% CUMCM 2026 B. Class = latexstudio/CUMCMThesis.
\documentclass{cumcmthesis}
\usepackage{unicode-math}
\setmathfont{Cambria Math}
\setmonofont{Courier New}
\setCJKmonofont{FangSong}
% Keep the original PDF's Chinese fonts consistent across operating systems.
\setCJKmainfont{SimSun}[BoldFont=SimHei,ItalicFont=KaiTi]
\setCJKsansfont{SimHei}[BoldFont=SimHei]
\setCJKfamilyfont{zhsong}{SimSun}[BoldFont=SimHei]
\setCJKfamilyfont{zhhei}{SimHei}[BoldFont=SimHei]
\setCJKfamilyfont{zhkai}{KaiTi}[BoldFont=KaiTi]
\setCJKfamilyfont{zhfs}{FangSong}[BoldFont=FangSong]
\usepackage{amsmath}
\usepackage{graphicx}
\usepackage{booktabs}
\usepackage{longtable}
\usepackage{array}
\usepackage{placeins}
\usepackage{xurl}

\title{无线电干扰源的交会可行集与机器狗清除策略}
\tihao{B}
\baominghao{202617201069}
\schoolname{Operator}
\membera{A}
\memberb{B}
\memberc{C}
\supervisor{—}

\begin{document}
\maketitle

\begin{abstract}
无线电测向把示向度误差限制在正负1度。本文将每条读数解释为半角1度的凸楔形，定位区域取诸楔之交；第二检测点沿第一示向度的正交基线布置；全向源采用9点侦听网给出最小接收半径下的距离可听闭式界，并对未扫静默频道的连续残留做剪枝，时间层先走覆盖骨架、一站示向度多数延后到覆盖结束后再沿测得示向度搜索；混编情形沿用该9点网并加外环侦听点，定向源听到当且仅当某网点落在距离盘且前瓣内，残留是开背瓣，听到后进入保扇形追逐。
问题一：合成正交交会例的对径圆覆盖为1；未与场地圆盘相交的两站有界样本覆盖率为1，相交裁剪后覆盖率为0.955；边长80 m的等边三角形对径圆覆盖为0。
问题二：第二检测点取轴上交角不小于45度的覆盖点，轴向距离800 m、横向偏移700 m。网格分位点横向偏移690.90909 m的事后直径中位数为48.2745 m，楔内70个样本上一次清除比例为0.2142857。固定轴向1100 m时，横向200 m的事后直径显著大于横向800 m。
问题三：9点侦听网的闭式覆盖半径为956.0510880488383 m，小于1000 m。
问题四：检测网沿用问题三的9点侦听网并加外环等分点。光学方格边长25 m的外接半径为17.677669529663685 m，小于20 m。
\keywords{无线电测向，交会定位，凸楔形，Jung定理，侦听网}
\end{abstract}

\section{问题重述}

\subsection{背景}

无线电干扰源的定位与清除是频谱管理中的紧迫任务。传统流程把工作分成初步监测、区域排查和精准定位三个阶段：前两个阶段用固定站、监测车或无人机圈出目标分布区域，最后一公里仍常由技术人员携带便携测向机徒步完成。环境遮挡、低功率辐射和定向波束都会使单点场强判断失效，完成时间又直接决定干扰的残留影响。题目因此提出：在半径1800 m的已知圆域内，用搭载测向机与激光清除装置的机器狗代替人工，建立可执行的自动定位与清除模型。

测向交会并不是新问题。Stansfield早年就把多站示向度的交会写成统计定位，并指出“三角形帽”（cocked hat）的大小本身并不等于定位可靠度\cite{stansfield1947}。无源测向的误差传播随后被系统整理：方位观测是非线性的，几何构型决定可观测性与协方差形状\cite{torrieri1984,foy1976}。单站纯方位跟踪还要求观测平台作侧向机动，否则距离不可辨\cite{nardone1984}。这些经典结论进入本题时必须改写：题目把误差明确限制在$[-1^\circ,1^\circ]$，并且要求输出的是定位区域的直径，而不是一个最大似然点。因此本文把示向度读数解释为楔形可行集的轴，把交会解释为凸集相交，把清除解释为把该凸集送进半径20 m的光学窗口。

\subsection{问题提出}

坐标系以场地圆心为原点，正东为$x$轴正向，正北为$y$轴正向，单位为米。示向度是从$x$轴逆时针转到“检测点指向干扰源”的角度，取值$[0^\circ,360^\circ)$。干扰源个数未知，频道取自$1$至$20$且互不相同、不随时间改变。全向源覆盖$360^\circ$；定向源覆盖其指向两侧各$90^\circ$。有效接收半径落在$[1000,1500]$ m，事先未知。机器狗速度5 m/s，换台1 s，检测5 s，光学精定位3 s，激光清除2 s。光学清除只依赖距离，不依赖波束朝向。同一地点重复检测不改变示向度误差。附件1与附件2给出HTTP接口与测试流程，并不提供问题1的站点坐标表。冻结附件的数据集摘要为\path{sha256:50bf04217eb3e8e7d3ce8b4a84f228b46dcfd93eebae187b443adb7724b2440e}。

四个问题彼此独立，不能并成同一目标函数。

\begin{enumerate}
\item 已知若干检测点及对应示向度，给出交会多边形直径的算法，并判断以该直径为直径的圆能否覆盖定位区域。
\item 已知全向源在一个检测点的示向度，给出第二个检测点的选择策略与候选区域。
\item 场地内10至16个全向源，一只机器狗从原点、频道1出发，自动定位清除；用演练统计检验清除比例与平均时间。
\item 混入未知朝向的定向源，其余条件同问题三，再给出策略并通过演练检验。
\end{enumerate}

题目要求在充分演练之后，按表1的五列格式报告三次正式测试：案例编码、清除干扰源个数、平均定位清除时间、程序运行时间。正式测试占用不可撤回的次数。本文把几何判定、侦听覆盖与演练窗口统计分开书写：演练行只出现在问题三、问题四的演练表中；正式测试表按模拟器完成弹窗的案例编码与HTTP导出的清除个数、平均定位清除时间、程序运行时间填写。正式完成弹窗不给出干扰源总数，因而该清除个数不得写成总数。论文体例遵循竞赛格式规范\cite{format}。

\section{问题分析}

四问对应三类任务：凸可行集的几何度量、第二站的正交基线设计、以及带时间代价的搜索清除。如图\ref{fig:framework}所示，测向读数先变成楔形，楔形相交得到定位区域，区域直径决定一次清除是否充分；第二站把交角抬到接近直角。侦听覆盖必须把距离可听与朝向可听分开：最小有效接收半径只约束距离因子$R_{\mathrm{eff}}$，并不约束波束朝向。在全向假设下，原点加8环的9点网对场地最近距离有闭式上界$\rho\approx 956.05\,\mathrm{m}<1000\,\mathrm{m}$，当场地点到某侦听点的距离不超过该半径时必被听到。定向源即使落在$R_{\mathrm{eff}}$球内，仍可因$180^\circ$瓣的背面而听不到。问题四因此在同一9盘之外加半径$2000\,\mathrm{m}$的12等分点，把检测写成混合源听力：全向完成9盘必听；定向听到当且仅当某网点落在距离盘且前瓣内。残留是开背瓣，不是对一切朝向的覆盖。全向源在覆盖骨架上扫完静默频道后再局部化；定向源在网点听到后进入保扇形追逐。四问在几何上依次衔接：全向距离覆盖在问题三关闭，定向瓣的混合听力留给问题四。

\begin{figure}[htbp]
\centering
\includegraphics[width=0.98\textwidth]{figures/framework.png}
\caption{四问几何衔接。(a) 示向楔（半角为示意放大，正文取$1^\circ$）；(b) 交会区与对径圆；(c) 交付点$(r,t)=(800,700)$ m，不是网格分位点$|t|=690.9$ m；(d) 问题四检测网$N_9\cup R_{12}$，虚线为1000 m，实线为1800 m，外环在2000 m。}
\label{fig:framework}
\end{figure}

\subsection{问题一的分析}

交会定位的教科书图像是两条射线交于一点。题目却把误差写成闭区间，并且强调同一地点的误差固定、不同地点才呈现统计规律。于是单次读数$\theta$不能当作真方位，只能当作楔形$W(S,\theta;1^\circ)$的轴。若干检测点对同一源的定位区域是这些楔形的交
\begin{equation}
\Omega=\bigcap_i W(S_i,\theta_i;1^\circ).
\label{eq:omega}
\end{equation}
两站时$\Omega$就是题附图中的红色四边形；多站时是至多$2n$边的凸多边形。直径
\begin{equation}
D=\mathrm{diam}(\Omega)=\sup_{P,Q\in\Omega}\|P-Q\|
\label{eq:diam}
\end{equation}
在紧凸多边形上由顶点对实现\cite{toussaint1983}。问题的后半句问：是否存在直径等于$D$的圆覆盖$\Omega$。这不是恒真命题。Jung定理给出平面集合最小包围圆半径的上界$R_{\mathrm{sec}}\le D/\sqrt{3}$，等边三角形取等号\cite{jung1901}。对径圆覆盖当且仅当$R_{\mathrm{sec}}\le D/2$，等边三角形破坏该不等式。因此算法必须同时输出$D$、最小包围圆半径以及覆盖布尔值，而不能预先回答“能”。

把交会写成无约束直线求交会得到后方虚假顶点。半平面裁剪只保留前方楔内的点：在正交例上与无约束交点重合，在对向示向度下则拒绝后方交点\cite{sutherland1974}。最小包围圆可用Welzl随机增量算法在期望线性时间内得到\cite{welzl1991}，多边形顶点很少，枚举直径对同样足够。

\subsection{问题二的分析}

第一站之后，可行集是一条窄而长的楔，再与距离环$5<\|P-S_1\|\le 1500$及场地圆盘相交。若第二站仍落在第一示向度轴上，两楔几乎平行，交会无界或极长。纯方位定位的经典可观测性分析已经指出：观测器必须离开视线才能分辨距离\cite{nardone1984}。最小二乘点估计还会随几何产生偏差\cite{dogancay2005}。因此第二站应写成
\begin{equation}
S_2=S_1+r\,u_{\theta}+t\,v_{\theta},
\label{eq:s2}
\end{equation}
其中$u_{\theta}$是示向度单位向量，$v_{\theta}$是左转法向，$t$可正可负。交角$\psi$在$r$等于源的轴向距离时达到$90^\circ$；$\psi\in[45^\circ,135^\circ]$当且仅当$|r-r_G|\le|t|$。小角近似
\begin{equation}
D\approx 2\varepsilon\sqrt{r_1^2+r_2^2}
\label{eq:smallang}
\end{equation}
在$r_1=1000$ m、$r_2=500$ m、$\varepsilon=1^\circ$时约为39 m，与清除直径40 m同量级，这是选取$|t|$的尺度。交付点由轴上$\psi\ge 45^\circ$覆盖引理给出：在$r=800$ m处$|t|\ge 700$ m。网格分位带
\[
r\in[800,1300]\,\mathrm{m},\qquad |t|\in[459,845]\,\mathrm{m}
\]
是同一套$n=70$个楔内样本上的比较，其中$(800,459)$已使远轴$\psi<45^\circ$，因而该带不是覆盖集。$[45^\circ,135^\circ]$覆盖另需$|t|\ge 795$ m。两侧对称，追逐时取离当前更近的一侧。

\subsection{问题三的分析}

全向源在覆盖角内场强只随距离衰减，因而最小接收半径在本问里只约束距离可听。原点不能单独听到全部源：场地半径1800 m，若有效接收半径取下限1000 m，远处源从原点必听不到。因此先要用侦听网在最坏接收半径下覆盖场地。覆盖模型是原点加8环共9盘的最近距离$d_{\mathcal{N}}(G)=\min_{S\in\mathcal{N}}\|G-S\|$在场地闭圆盘上的最大值$\rho$。该最大值的候选类是原点--环点Voronoi顶点与场地边界角平分点$P$：原点与相邻两环点的Voronoi顶点$V$落在角平分线上，半径$\tau=1000/(2\cos(\pi/8))\approx 541.196$ m；$P$落在场地边界同一角平分处，到最近环点的距离即为$\rho\approx 956.05$ m。分款以$\tau$为界：$r\le\tau$时最近侦听点是原点；$r\in(\tau,1800]$时最近侦听点是环点，其中$r\in(\tau,1000]$仍由环点覆盖而不是原点。1000 m圆是环点所在的设计听距圆，不是$\rho$的分款。仅8个环点不能覆盖原点，原点到环点的距离是1000 m。抽样比例只作核对。该闭式界只在全向假设下成立，不能搬到问题四的定向瓣上。

发现层的跳过规则是连续残留而不是有限核对点。记发现站为实际扫描过剩余静默频道的位姿，连续残留是场地闭圆盘减去这些站的闭圆盘。侦听点$Q\in\mathbb{A}$仅当已证明上界$U=\rho(\mathcal{S},K_Q)\le 1000$ m时才丢掉；数值不确定则保留。有限边界样本不能代替该残留。场地边界角平分点$P=1800\angle 22.5^\circ$到$R_0$与$R_1$的距离都等于覆盖半径$\rho\approx 956.05$ m，故访问$R_1$即可使该有限点离开残留；同一边界上的$G_\star=1800\angle 7.5^\circ$只落在$R_0$的1000 m盘内，到$R_1$约1176 m。若因这些角平分样本已空而跳过$R_0$，则$G_\star$不会被听到。八个角平分点$1800\angle(22.5^\circ+k\cdot 45^\circ)$在去掉$R_0$后仍全部落入其余网点的1000 m盘，但$G_\star$到其余8个网点的最近距离仍大于1000 m。覆盖目标从当前位姿对剩余覆盖路点做开放TSP。若已选出共享D-最优站，记该站相对TSP首站的路程差为到该站的距离减去到首站的距离：当路程差不超过80 m时，下一站抢先取该D-最优站；当该站落在剩余覆盖路点120 m邻域内或覆盖某个剩余透镜，并且路程差不超过220 m时，同样抢先取该站。共享D-最优站、边上顺路探测与最近剩余是时间启发式。一站示向度多数延后到覆盖骨架走完之后；覆盖未完时，仅当接合点距当前不超过420 m，或第一测向站到原点的距离大于200 m且接合点连续割完剩余透镜时，才并入当前路径。两站及以上的服务点在覆盖未完时进入当前开放TSP，不要求定位多边形已达清除证书；选中后先圆心再测，仍开放则沿测得示向度归航。覆盖未完且启发式失败时，覆盖义务仍把与残留相交的9网点送回覆盖目标，不把覆盖骨架改写成命中即射线的星形。覆盖循环另设动作闸$A_{\max}=900$：该上限是算法自设，原题未规定。达到上限、程序运行时间、测试窗口或技术异常时，覆盖扫描、静默扫描、射线归航与追逐一并返回。这些情形是检测全称的残差。覆盖循环若按残留空、静默频道空或已听满16个正常结束，则每个全向源已被听到。本文不证明覆盖义务必在有限步内走完。

发现与清除分成两层。覆盖路点对静默频道做检测时，\texttt{near}就地清除，\texttt{direction}只记示向度，该点把剩余静默频道扫完后再走下一覆盖路点。一站示向度多数延后到覆盖骨架走完之后的服务旅行商路径；覆盖未完时仅当接合点距当前不超过420 m，或第一测向站到原点大于200 m且接合点连续割完剩余透镜时并入当前路径。两站及以上的服务点在覆盖未完时进入当前开放TSP，不要求已达清除证书；选中该服务点后先按圆心再测处理定位多边形，仍开放则沿测得示向度归航。清除证书仍是$R_{\mathrm{sec}}\le 20$ m（或对径圆覆盖且$D/2\le 20$ m），未达证书并不阻止两站服务点进入开放TSP。已达证书的圆心若落在近邻，另可顺路清除。轴上清除点$P$满足$\|SP\|=\|SG\|=r$时，弦长$\|PG\|=2r\sin(|\alpha|/2)$：对某个实际源，命中当且仅当该弦长不超过20 m；对一切$|\alpha|\le 1^\circ$，保证命中当且仅当$r\le 20/(2\sin 0.5^\circ)\approx 1145.930\,\mathrm{m}$。$r\tan 1^\circ$是轴向距离$r$处垂直截面半宽，只给出保守充分条件，其数值$20/\tan 1^\circ\approx 1145.799\,\mathrm{m}$。9盘只保证存在某个网点到源的距离不超过$\rho$，这不是首次听到的站。例如$G=(1400,0)$、$R_{\mathrm{eff}}=1500$ m时原点首次听到，距离1400 m，$2800\sin 0.5^\circ\approx 24.43>20$。存在一个原始网点，在真实距离已知的理想条件下该轴上清除圆能盖住源；真实距离未知，这仍不是可执行的清除保证。可执行的清除保证仍是圆心证书。沿射线的扩展搜索、横向有限振幅、最后命中点\texttt{/clear}与沿示向度15 m都不是这个弦圆。问题二的交付点$(800,700)$ m不是本问默认第二站。虚拟时间只在接口接受的检测与清除上推进，移动代价按直线距离除以5 m/s计算。现实时间由\texttt{/enter}返回的剩余秒数约束，虚拟时间可以远长于现实时间。

\subsection{问题四的分析}

定向源的有效覆盖是未知朝向的闭$180^\circ$扇形。最小有效接收半径在本问里仍然只约束距离因子：检测点收不到信号，既可能是超距，也可能落在扇形背面，因此\texttt{no\_signal}不能把该频道标死，也不能把问题三在正常结束时的全向检测搬到定向瓣上。光学清除仍只依赖距离，问题一的圆心清除规则对定向源仍然充分。

本问只采用一套检测网
\[
\mathcal{N}=N_9\cup R_{12},
\]
其中$N_9$是问题三的原点加半径$1000\,\mathrm{m}$的8等分环，$R_{12}$是半径$2000\,\mathrm{m}$上的12等分点，圆心角$30^\circ$，最近外环点相对源的方位差$\Delta\theta\le 15^\circ$。全向源走完$N_9$必听，当且仅当覆盖半径$\rho(N_9)\approx 956.05\,\mathrm{m}<1000\,\mathrm{m}$。定向源听到当且仅当存在$S\in\mathcal{N}$使$\|S-g\|\le R_{\mathrm{eff}}$且$(S-g)\cdot u\ge 0$。边界中缝到最近外环点的距离最大约$534.17\,\mathrm{m}<1000\,\mathrm{m}$，故每个边界位置都有外环点落在距离盘内；若该最近点还落在前瓣（含径向外指），则该定向源被听到。残留是
\[
\mathrm{Res}=\bigl\{(g,u):\mathcal{N}\cap\overline{\mathbb{D}}(g,1000)\ \text{全落在开背瓣}\bigr\}.
\]
这不是$\forall g\forall u$。占用达到16则剩余静默频道为空。单次\texttt{no\_signal}不能清空频道。三次演练全部清除不能代替该检测引理。

时间层把发现与追逐分开：内圈按$N_9$的连续残留做蛇形访问，仅附近才打断去追逐；外环$R_{12}$开圆巡回；覆盖结束后一条服务旅行商路径。追逐仍用短沿迹加侧向的保扇形站：内点计划的前两项为$(r,t)=(150,\pm 400)$ m，模长$\sqrt{150^2+400^2}\approx 427.2$ m。沿示向度走过源会落入背面，因而本问第二站不取问题二的交付点$(800,700)$ m。两站非共线时先做与问题三相同的多边形半径判定：$20\,\mathrm{m}<R_{\mathrm{sec}}\le 1000$ m仍在最小包围圆圆心测量；圆心返回\texttt{no\_signal}且$R_{\mathrm{sec}}\le 80$ m时，先在圆心强制\texttt{/clear}，再走18 m六邻点，然后才把判定标成开放，并对最后一站再取保扇形站。循环用尽后先在圆心强制\texttt{/clear}；若$R_{\mathrm{sec}}\le 80$ m仍未清除，则铺边长25 m的光学方格，随后仍沿测得示向度按路程$\{120,\ldots,1100\}$ m作瓣内最后手段。混编覆盖循环另设动作闸$A_{\max}=900$：该上限是算法自设，原题未规定。协议允许场地外检测，外环因此可以落在圆盘之外。

\section{模型假设}

\begin{enumerate}
\item 同一检测点、同一电磁环境下，示向度误差不随重复检测改变；不同地点的误差各自落在$[-1^\circ,1^\circ]$。定位区域按该闭区间的最坏楔形构造，而不是按单次随机实现。
\item 测向机读数$\theta$表示从检测点指向干扰源的方位轴。真实方位与读数的圆距离不超过$1^\circ$，故可行集是半角$1^\circ$的闭凸楔。
\item 干扰源位于场地闭圆盘内。已测到信号时，源与检测点距离不超过1500 m。
\item $\Omega$是凸集。紧凸多边形的直径在顶点对上取得。
\item 以直径$D$为直径的圆覆盖$\Omega$，当且仅当最小包围圆半径$R_{\mathrm{sec}}\le D/2$。
\item 全向源在360度内均匀辐射；定向源在朝向两侧各90度内均匀辐射，背面无信号。场强在覆盖角内只随距离衰减。
\item 光学精定位与激光清除的成败只取决于清除点与源的距离是否不超过20 m，与波束朝向无关。距离不超过5 m且位于覆盖角内时无法获得示向度，可直接清除。
\item 有效接收半径在$[1000,1500]$ m内未知。侦听网按1000 m设计，第二站可听约束同时报告1500 m上界。
\item 机器狗沿直线行进，速度5 m/s；换台、检测、光学与激光时间取题目给定常数。虚拟时间只在被接受的检测与清除指令上累加。
\item 问题一、问题二的数值直径使用文中写明的合成站点，不把这些坐标当作赛题真值。附件不提供问题一坐标表。
\item 问题三、问题四的接口动作次数以外层闸$A_{\max}=900$截断。该上限是算法自设，原题未规定。覆盖扫描、静默频道检测、射线归航与追逐在达到该上限时返回；程序运行时间、测试窗口与技术异常同样中断。检测全称把这些情形列为残差，并且不证明覆盖循环必在有限步内走完。
\end{enumerate}

\clearpage
\section{符号说明}

几何与运动学符号如下。

\begin{center}
\begin{tabular}{@{}>{\raggedright\arraybackslash}p{0.26\textwidth}>{\raggedright\arraybackslash}p{0.54\textwidth}p{0.10\textwidth}@{}}
\toprule
符号 & 含义 & 单位 \\
\midrule
$S_i=(x_i,y_i)$ & 第$i$个检测点 & m \\
$\theta_i$ & 在$S_i$测得的示向度 & 度 \\
$\varepsilon$ & 示向度半误差，题设为$1$ & 度 \\
$W(S,\theta;\varepsilon)$ & 以$S$为顶点、轴$\theta$、半角$\varepsilon$的闭楔形 & --- \\
$n_\pm$ & 楔左右边界指向内部的单位法向 & --- \\
$\Omega$ & 定位区域 $\bigcap_i W(S_i,\theta_i;\varepsilon)$ & --- \\
$D$ & $\mathrm{diam}(\Omega)$ & m \\
$R_{\mathrm{sec}}$ & $\Omega$的最小包围圆半径 & m \\
$R_A$ & 场地半径 & m \\
$R_{\mathrm{eff}}$ & 有效接收半径 & m \\
$R_{\mathrm{clr}}$ & 清除半径 & m \\
$R_{\mathrm{near}}$ & 近场无示向度阈值 & m \\
$v$ & 机器狗速度 & m/s \\
$T_{\mathrm{sw}},T_{\mathrm{det}}$ & 换台时间、检测时间 & s \\
$T_{\mathrm{opt}},T_{\mathrm{las}}$ & 光学精定位、激光清除 & s \\
$\psi$ & 两站中心线在目标处的交角 & 度 \\
$r,t$ & 第二站轴向距离、横向偏移 & m \\
$u_{\theta},v_{\theta}$ & 示向度单位向量及其左转法向 & --- \\
\bottomrule
\end{tabular}
\end{center}

侦听、覆盖与清除符号如下。

\begin{center}
\begin{tabular}{@{}>{\raggedright\arraybackslash}p{0.26\textwidth}>{\raggedright\arraybackslash}p{0.54\textwidth}p{0.10\textwidth}@{}}
\toprule
符号 & 含义 & 单位 \\
\midrule
$c$ & 频道编号 & --- \\
$N$ & 干扰源个数 & --- \\
$H$ & 定向源朝向 & 度 \\
$N_{\mathrm{scout}}$ & 侦听点个数 & --- \\
$N_9$ & 问题三原点加8环的9点侦听网 & --- \\
$R_{12}$ & 半径2000 m的12等分外环 & --- \\
$\mathcal{N}$ & 问题四检测网 $N_9\cup R_{12}$ & --- \\
$\rho$ & 问题三9盘最近距离上界 & m \\
$d_{\mathrm{mid}}$ & 外环边界中缝到最近点的距离 & m \\
$\Delta\theta$ & 源与最近外环点的方位差，至多$15^\circ$ & 度 \\
$S^\star$ & $R_{12}$中方位最近的点 & m \\
$\mathrm{Res}$ & 定向残留：网点与距离盘之交全在开背瓣 & --- \\
$\mathrm{Res}_{\mathrm{cont}}$ & 连续残留：场地减去发现站闭圆盘 & --- \\
$K_Q$ & 侦听点$Q$的透镜 $\overline{\mathbb{D}}(Q,1000)\cap\mathbb{A}$ & --- \\
$\gamma$ & $\cos(\pi/8)$，9盘覆盖证明中的角平分余弦 & --- \\
$\tau$ & 原点与环点在角平分线上的Voronoi顶点半径 $1000/(2\gamma)$ & m \\
$r_{\mathrm{ray}}$ & 轴上弦长清除半径界：$P$满足$\|SP\|=\|SG\|=r$时，对一切$|\alpha|\le 1^\circ$保证命中当且仅当$r\le 20/(2\sin 0.5^\circ)$ & m \\
$\Delta_{\mathrm{edge}}$ & 边上顺路插入的额外路程 & m \\
$I(p)$ & 共享第二站对若干$2^\circ$楔的期望对数行列式 & --- \\
$u_H$ & 定向源朝向单位向量 & --- \\
$\eta_1$ & 一次清除指示：$D\le 40$ m且对径圆覆盖，或$R_{\mathrm{sec}}\le 20$ m & --- \\
$R_{\mathrm{hear}}$ & 圆心再测半径，取最小有效接收半径1000 & m \\
$R_{\mathrm{hex}}$ & 六邻点尝试半径80，不是80 m盘的覆盖半径 & m \\
$N_{\mathrm{true}},N_{\mathrm{clr}}$ & 干扰源总数、被清除个数 & --- \\
$A,A_{\max}$ & 已发出的接口动作次数及其上限，上限取900 & --- \\
$p_{\mathrm{hear}}$ & 第二站对楔内样本的可听比例 & --- \\
\bottomrule
\end{tabular}
\end{center}

\section{数据预处理}

\subsection{附件口径与字段}

本题没有问题一的站点坐标表，也没有可直接读取的源位置。附件1、附件2给出模拟器的通信协议、动作语义与测试窗口，属于接口说明而不是观测样本。预处理的任务因此是把题面常数、接口时间模型和几何约束整理成可复用的参数表，并冻结数据摘要，而不是做缺失插补。后续定量结果均在该冻结附件上复现。\footnote{数据集摘要哈希为\path{sha256:50bf04217eb3e8e7d3ce8b4a84f228b46dcfd93eebae187b443adb7724b2440e}。}

表\ref{tab:protocol}列出从题面与附录抽取的常数。有效接收半径是区间而不是单值，侦听与第二站设计必须在区间两端都可解释。清除半径与近场阈值是硬约束：前者决定一次\texttt{/clear}的成功圆盘，后者决定何时放弃示向度、直接光学清除。

\begin{table}[htbp]
\centering
\caption{题面与附录抽取的协议常数}
\label{tab:protocol}
\begin{tabular}{lrrl}
\toprule
量 & 取值 & 单位 & 来源 \\
\midrule
场地半径 $R_A$ & 1800 & m & 题面 \\
示向度半误差 $\varepsilon$ & 1 & 度 & 附录2 \\
有效接收半径 & $[1000,1500]$ & m & 附录2 \\
清除半径 $R_{\mathrm{clr}}$ & 20 & m & 附录2 \\
近场阈值 $R_{\mathrm{near}}$ & 5 & m & 附录2 \\
行进速度 $v$ & 5 & m/s & 附录2 \\
换台时间 & 1 & s & 附录2 \\
检测时间 & 5 & s & 附录2 \\
光学精定位 & 3 & s & 附录2 \\
激光清除 & 2 & s & 附录2 \\
频道数 & 20 & --- & 附录1 \\
源个数 & $[10,16]$ & --- & 问题三、四 \\
坐标分量绝对值上限 & $2\times 10^6$ & m & 附录3 \\
虚拟时间上限 & 100 & h & 附录3 \\
\bottomrule
\end{tabular}
\end{table}

\subsection{尺度一致性与质量控制}

如图\ref{fig:radii}所示，五类半径画在同一圆心上：近场5 m和清除20 m相对于场地几乎是一个点。这表明搜索阶段的主要矛盾是“听得到”，局部化阶段的主要矛盾才是“圆得足够小”。如图\ref{fig:constline}所示，同一组半径排在对数坐标上，尺度链
\[
R_{\mathrm{near}}<R_{\mathrm{clr}}<R_{\mathrm{eff}}^{\min}\le R_{\mathrm{eff}}^{\max}<R_A,
\]
即$5<20<1000\le 1500<1800$依次成立。该链在题面中全部给定，预处理只做完备性核对，不改写任一端点。

\begin{figure}[htbp]
\centering
\includegraphics[width=0.72\textwidth]{figures/eda_radii.png}
\caption{场地、接收、清除与近场半径的相对尺度。}
\label{fig:radii}
\end{figure}

\begin{figure}[htbp]
\centering
\includegraphics[width=0.92\textwidth]{figures/eda_constraint_line.png}
\caption{协议半径的对数尺度链：近场、清除、接收区间与场地半径依次增大。}
\label{fig:constline}
\end{figure}

如图\ref{fig:timecomp}所示，一次典型动作的虚拟时间中，行进200 m需要40 s，检测5 s，光学3 s，激光2 s，换台1 s。这表明移动项通常主导虚拟时钟，侦听点不宜过密，追逐时应沿用已有示向度。如图\ref{fig:protorange}所示，长度量在对数坐标下拉开：近场与清除远小于接收半径，场地半径最大。预处理不把接收区间用中点代替，而是在覆盖论证中取1000 m，在可听上界中取1500 m。

\begin{figure}[htbp]
\centering
\includegraphics[width=0.82\textwidth]{figures/eda_time_components.png}
\caption{虚拟时间分量：行进200 m、检测、光学、激光与换台。}
\label{fig:timecomp}
\end{figure}

\begin{figure}[htbp]
\centering
\includegraphics[width=0.72\textwidth]{figures/eda_protocol_ranges.png}
\caption{协议长度量的对数对照：仅有效接收半径为区间。}
\label{fig:protorange}
\end{figure}

\subsection{接口语义与标签口径}

模拟器通过HTTP返回\texttt{direction}、\texttt{near}、\texttt{no\_signal}以及清除结果。表\ref{tab:api}归纳四类指令的时钟含义。\texttt{direction}只提供楔形轴，不提供距离。\texttt{near}表示距离不超过5 m且在覆盖角内。\texttt{no\_signal}在全向情形可解释为超距或未进入该频道；在定向情形还可能是扇形背面。清除成功与否不能从示向度反推，只能由\texttt{/clear}的返回判定。问题三、问题四的“被清除个数比例”和“平均定位清除时间”是模拟器窗口上的统计，不是附件中的既有列。演练结束时窗口给出干扰源总数，因而比例可算；正式测试不给出总数，表1只报告清除个数与平均时间。

\begin{table}[htbp]
\centering
\caption{接口指令与虚拟时钟}
\label{tab:api}
\begin{tabular}{lll}
\toprule
指令 & 作用 & 虚拟时间增量 \\
\midrule
\texttt{/enter} & 进入场地，返回剩余现实秒数 & 0 \\
\texttt{/measure} & 在停稳点检测一个频道 & $\|P_{\mathrm{new}}-P_{\mathrm{old}}\|/5+1_{\mathrm{switch}}+5$ \\
\texttt{/clear} & 光学精定位并尝试激光清除 & $\|P_{\mathrm{new}}-P_{\mathrm{old}}\|/5+$（成功5或失败3） \\
\texttt{/exit} & 结束本次测试 & 0 \\
\bottomrule
\end{tabular}
\end{table}

每次动作使用新的请求编号，串行等待上一次返回后再发下一次；仅在重试同一中断请求时复用编号。坐标分量绝对值不得超过$2\times 10^6$ m，虚拟时间上限为100 h。实现细节见附录。

\section{问题一：交会楔形的直径与对径圆覆盖}

\subsection{机理}

测向机在检测点$S$把天线转到场强最大方向，读数$\theta$是该方向相对正东的方位。环境与仪器使读数相对真方位有不超过$1^\circ$的偏差，且同一地点的偏差固定。本文把该读数解释为凸楔形（convex wedge）的轴。其基本形式为
\begin{equation}
W(S,\theta;\varepsilon)=\bigl\{P:\delta\bigl(\tfrac{180}{\pi}\mathrm{atan2}(P_y-S_y,P_x-S_x),\,\theta\bigr)\le\varepsilon\bigr\},
\label{eq:wedge}
\end{equation}
其中圆距离$\delta(\alpha,\theta)=\min_k|\alpha-\theta-360k|$，$\varepsilon=1^\circ$为题设半误差，$P$为平面上的候选点。半角$1^\circ$远小于$90^\circ$，楔是凸的，多楔之交仍凸。这与把示向度当无误差射线、再对交点做泰勒线性化的点估计路线不同\cite{foy1976,torrieri1984}：本题要的是区域直径，不是协方差椭圆的一次实现。如图\ref{fig:halfplane}所示，轴为$\theta$，左右边界的内法向为$n_+$与$n_-$：可行点必须同时落在两条边界直线的内侧半平面。半角在图中放大以便阅读，正文始终取$1^\circ$。

\begin{figure}[htbp]
\centering
\includegraphics[width=0.78\textwidth]{figures/halfplane.png}
\caption{单站示向度楔：轴为$\theta$，左右边界的内法向为$n_+$与$n_-$。}
\label{fig:halfplane}
\end{figure}

两站中心线不平行时，四条边界射线在前方相交，得到有界四边形。如图\ref{fig:wedgeint}所示，左幅是$1^\circ$真半角的两站楔，交会区在本图尺度下缩为一点，右幅把该点放大成定位四边形。1500 m距离圆远大于该视窗，正交例的交会区未与距离圆相交。三站及以上可以把$\Omega$削成锐角三角形或六边形。锐角三角形的外接圆半径大于最长边的一半，这正是Jung等号的几何来源\cite{jung1901}。如图\ref{fig:jungschema}所示，左侧最小包围圆明显大于以一边为直径的圆，右侧对径圆恰好包住钻石。这表明对径圆覆盖不是恒真命题。

\begin{figure}[htbp]
\centering
\includegraphics[width=0.82\textwidth]{figures/wedge_intersect.png}
\caption{两站前方楔形之交构成定位区域$\Omega$。左：真半角$1^\circ$；右：交会区放大。}
\label{fig:wedgeint}
\end{figure}

\begin{figure}[htbp]
\centering
\includegraphics[width=0.82\textwidth]{figures/jung_vs_diamond.png}
\caption{对径圆覆盖不是恒真：等边三角形失败，钻石四边形成功。}
\label{fig:jungschema}
\end{figure}

\subsection{模型}

楔的半平面写法便于裁剪。记左边界方位$\phi_+=\theta+\varepsilon$，右边界方位$\phi_-=\theta-\varepsilon$，对应单位方向$u=(\cos\phi,\sin\phi)$。指向楔内的法向为
\begin{equation}
n_+=(\sin\phi_+,-\cos\phi_+),\qquad n_-=(-\sin\phi_-,\cos\phi_-).
\label{eq:normals}
\end{equation}
点$P$落在楔内当且仅当$n_\pm\cdot(P-S)\ge 0$。定位多边形由Sutherland--Hodgman半平面裁剪（polygon clipping）得到。其基本形式为：从足够大的正方形出发，依次用半平面$n\cdot x\ge c$裁去外侧顶点并插入边界交点，输出仍为凸多边形\cite{sutherland1974}。若顶点落在人工大盒边界上，则判定$\Omega$无界，有限直径圆不能覆盖。干扰源在场地内时再交圆盘$\mathbb{D}(0,1800)$；已测到信号时再交半径1500 m的距离球。圆盘与距离球用切线半平面多边形逼近后同样裁剪。

算法步骤如下。

\begin{enumerate}
\item 读入检测点与示向度$\{(S_i,\theta_i)\}$，对每一站构造两条半平面。
\item 自大正方形起做半平面裁剪，得到凸多边形顶点；必要时再交场地圆盘与距离球。
\item 若顶点落在人工大盒边界上，判定无界，直径为$+\infty$，有限对径圆不能覆盖。
\item 否则计算
\begin{equation}
D=\max_{i<j}\|v_i-v_j\|.
\label{eq:dvert}
\end{equation}
\item 对每一对径点对$\{A,B\}$（$\|A-B\|=D$），检验其余顶点是否落入以$AB$为直径的闭圆盘。由Thales定理，这等价于$\angle APB\ge 90^\circ$。
\item 等价地，用Welzl算法计算最小包围圆（smallest enclosing circle）半径$R_{\mathrm{sec}}$\cite{welzl1991}。其基本形式为在顶点集上维护当前圆，遇圆外点则把该点固定到边界并递归，覆盖当且仅当
\begin{equation}
2R_{\mathrm{sec}}\le D.
\label{eq:cover}
\end{equation}
\end{enumerate}

Jung上界$R_{\mathrm{sec}}\le D/\sqrt{3}$只说明覆盖可能失败，并不给出失败的充分条件；失败必须由\eqref{eq:cover}的直接检验确认。无约束直线求交若不检查射线参数，会在对向示向度下给出后方交点。后方交点不满足式\eqref{eq:wedge}，不属于定位区域。本文以半平面裁剪为准。

\subsection{估计与检验}

题面未给检测点坐标。检验用文中写明的合成站点，种子为2026。正交交会例取$S_1=(0,0)$示向度$45^\circ$，$S_2=(1000,0)$示向度$135^\circ$，真源在$(500,500)$。等边反例取边长80 m的三角形；三站楔在等边核周围交出六边形。两站Monte Carlo在场地内抽样源与站点，叠加$\pm 1^\circ$读数后计算有界样本的覆盖率：请求400个样本，全部有界，其中382个对径圆覆盖，覆盖率$382/400=0.955$。直方图、散点与附录CSV使用这同一组400行，不是另行抽取的子样本。辅助直线交点与半平面顶点的最大距离用于核对正交例。

\subsection{结果}

由表\ref{tab:q1head}可知，正交交会例得到4个顶点，直径34.92 m，最小包围圆半径17.46 m，恰为直径的一半，对径圆覆盖取1，面积为609.36 m$^2$。辅助直线交点与半平面顶点的最大误差为0，说明在正交交叉这种“前方可见”的构型上，两种求交一致。如图\ref{fig:q1ex}所示，左幅给出两站与真源的相对位置，交会区在千米尺度下缩为一点，故用红框标出右幅视窗；右幅把定位四边形和对径圆放大到四十米视窗：四边形几乎是对径圆的内接钻石。这表明覆盖判定与目视一致。表\ref{tab:q1vert}列出四个顶点坐标，直径在相对的顶点对上取得。

\begin{table}[htbp]
\centering
\caption{问题一合成正交交会例的几何量}
\label{tab:q1head}
\begin{tabular}{lr}
\toprule
量 & 取值 \\
\midrule
顶点个数 & 4 \\
直径 / m & 34.92 \\
最小包围圆半径 / m & 17.46 \\
对径圆覆盖 & 1 \\
面积 / m$^2$ & 609.36 \\
辅助顶点最大误差 / m & 0 \\
比值 $R_{\mathrm{sec}}/D$ & 0.500 \\
\bottomrule
\end{tabular}
\end{table}

\begin{table}[htbp]
\centering
\caption{正交交会例的定位多边形顶点}
\label{tab:q1vert}
\begin{tabular}{crr}
\toprule
编号 & $x$ / m & $y$ / m \\
\midrule
1 & 500.000 & 517.765 \\
2 & 482.550 & 499.695 \\
3 & 500.000 & 482.844 \\
4 & 517.450 & 499.695 \\
\bottomrule
\end{tabular}
\end{table}

\begin{figure}[htbp]
\centering
\includegraphics[width=0.98\textwidth]{figures/q1_example_polygon.png}
\caption{正交交会合成例。左：两站与真源；右：定位四边形及其对径圆。}
\label{fig:q1ex}
\end{figure}

由表\ref{tab:q1cover}可知，未裁剪的两站有界样本384个全部被对径圆覆盖，覆盖率为1；与场地及距离球相交后，长条钻石被切断成略锐的三角形，400个有界样本中382个覆盖，覆盖率为0.955。边长80 m的等边三角形最小包围圆半径46.19 m，大于40 m，覆盖取0。三站楔的六边形同样覆盖取0，直径93.32 m，最小包围圆半径47.13 m。如图\ref{fig:jungschema}所示，左侧虚线对径圆不能包住等边三角形，右侧对径圆包住钻石。这表明“以直径为直径的圆能否覆盖”必须作为布尔输出，而不能预先回答“能”。

\begin{table}[htbp]
\centering
\caption{对径圆覆盖的正例与反例}
\label{tab:q1cover}
\begin{tabular}{lrr}
\toprule
情形 & 覆盖 & 说明 \\
\midrule
两站未裁剪有界样本 & 1 & 384个有界样本 \\
两站裁剪后 & 0.955 & 400个有界样本，其中382个覆盖 \\
等边三角形 & 0 & 边长80 m \\
三站六边形 & 0 & 反例 \\
\bottomrule
\end{tabular}
\end{table}

如图\ref{fig:q1mc}所示，裁剪后两站直径的直方图来自上述400个有界样本：中位直径137.95 m，显著大于清除直径40 m；右幅覆盖382个、未覆盖18个，与覆盖率0.955一致。这表明两站交会通常还不能一次清除，需要问题二把第二站放到近乎正交的位置，或者在追逐中继续加密测点。如图\ref{fig:jungratio}所示，左幅给出残差$R_{\mathrm{sec}}-D/2$：覆盖样本落在零线，18个未覆盖样本的残差为正，最大比值$R_{\mathrm{sec}}/D=0.5005$；右幅给出面积随直径的散点，长条区域面积并不随直径线性增长。

\begin{figure}[htbp]
\centering
\includegraphics[width=0.92\textwidth]{figures/q1_mc_diameter.png}
\caption{两站裁剪后400个有界样本的直径分布与覆盖计数（382覆盖、18未覆盖）。}
\label{fig:q1mc}
\end{figure}

\begin{figure}[htbp]
\centering
\includegraphics[width=0.92\textwidth]{figures/q1_jung_ratio_area.png}
\caption{左：400个有界样本的残差$R_{\mathrm{sec}}-D/2$；右：定位多边形面积与直径。}
\label{fig:jungratio}
\end{figure}

\subsection{诊断}

第一，直径算法在无界楔交上必须返回无穷，否则会把大盒边长误当作场地尺度。第二，辅助直线交点在正交例上与裁剪顶点误差为0，不能推广到对向示向度；对向时后方交点必须丢弃。第三，覆盖判定对顶点扰动稳定，因为正交例的$R_{\mathrm{sec}}/D$几乎恰好为$1/2$。第四，题目问“能否覆盖”的语言容易被读成恒真，等边反例说明必须保留布尔输出。第五，合成坐标不是赛题真值，算法对一般站点组可执行，数值只对写明的例子成立。第六，场地裁剪会把覆盖率从1降到0.955，正文报告覆盖结论时必须声明是否与场地相交。

\section{问题二：第二检测点的正交基线与候选带}

\subsection{机理}

第一站把源限制在窄楔$\Omega_1$内。楔的宽度在距离$r_G$处约为$2r_G\sin\varepsilon\approx 2r_G\varepsilon$，沿轴却可延伸到接收半径上界。第二站的任务是把这条长带切成直径可与清除圆相比的有界多边形。若$S_2$落在轴上，交角趋于$0^\circ$或$180^\circ$，切出来的仍是长条。Nardone等对纯方位可观测性的分析表明，观测器必须离开视线\cite{nardone1984}；Dogancay指出线性最小二乘在纯方位上有几何偏差\cite{dogancay2005}。因此把第二站写成正交基线（orthogonal baseline）。其基本形式为$S_2=S_1+r u_{\theta}+t v_{\theta}$，其中$u_{\theta}$为示向度单位向量，$v_{\theta}$为左转法向，$t$可正可负。横向偏移$t$是决策变量，而不是可有可无的扰动。

对轴上假设点$G=S_1+r_G u_{\theta}$，第二站$S_2=S_1+r u_{\theta}+t v_{\theta}$的交角满足
\begin{equation}
\cos\psi=\frac{r_G-r}{\sqrt{(r_G-r)^2+t^2}}.
\label{eq:psi}
\end{equation}
$\psi=90^\circ$当且仅当$r=r_G$。$\psi\in[45^\circ,135^\circ]$当且仅当$|r-r_G|\le|t|$。小角近似\eqref{eq:smallang}给出与40 m清除直径同量级的横向尺度，避免把$t$取得过小。

\subsection{模型}

如图\ref{fig:s2schema}所示，第二站写成$S_2=S_1+r u_{\theta}+t v_{\theta}$，候选比较在轴两侧进行，而不是绕$S_1$的圆环。对轴上$r_G\in[r_{\mathrm{lo}},r_{\mathrm{hi}}]$，交角覆盖写成两条引理。

\textbf{引理（轴上$[45^\circ,135^\circ]$覆盖）。} 假设$G$落在示向度轴上且$r_G\in[r_{\mathrm{lo}},r_{\mathrm{hi}}]$。则对该区间内一切$r_G$都有$\psi\in[45^\circ,135^\circ]$，当且仅当
\begin{equation}
|t|\ge\max\bigl(|r-r_{\mathrm{lo}}|,\,|r_{\mathrm{hi}}-r|\bigr).
\label{eq:psiband}
\end{equation}
取$r_{\mathrm{lo}}=5$ m、$r_{\mathrm{hi}}=1500$ m、$r=800$ m，得$|t|\ge 795$ m。

\textbf{引理（轴上$\psi\ge 45^\circ$覆盖）。} 同一假设下，$r_G\le r$时$\psi\ge 90^\circ\ge 45^\circ$自动成立。故对该区间内一切$r_G$都有$\psi\ge 45^\circ$，当且仅当
\begin{equation}
|t|\ge\max(0,\,r_{\mathrm{hi}}-r).
\label{eq:psi45}
\end{equation}
在$r=800$ m得$|t|\ge 700$ m。取$|t|=700$ m时，$\psi(1500)=45^\circ$，$\psi(5)\approx 138.64^\circ>135^\circ$。近端瘦钝角是本引理的残差，不是$[45^\circ,135^\circ]$覆盖。离轴$\pm 1^\circ$不在这两条引理范围内。

交付点取轴上$\psi\ge 45^\circ$覆盖点$(r,|t|)=(800,700)$ m。网格取$r\in[400,1400]$、$|t|\in[150,1000]$，对$\Omega_1$的径向--角度样本计算事后直径、交角和可听比例；该网格只作比较。网格分位点$(800,690.9)$ m使远轴$\psi(1500)\approx 44.63^\circ<45^\circ$，因而不是覆盖点。过样本形心的垂线把中位直径压得更小，但其直径90\%分位超过90 m，且不覆盖全轴，同样不是方案。两侧中选择离当前机器狗更近、更不远离场地的一点。若第二站\texttt{no\_signal}，改测另一侧$-t$，再沿示向度前进，而不是沿轴退回。问题二的交付点用于轴上交角覆盖；问题三全向局部化不把该点当作默认第二站。

\begin{figure}[htbp]
\centering
\includegraphics[width=0.72\textwidth]{figures/second_station.png}
\caption{第二站离开第一示向度轴：轴向距离$r$、横向偏移$t$。交付点$|t|=700$ m由轴上$\psi\ge 45^\circ$覆盖给出；网格分位点$|t|=690.9$ m只作比较。}
\label{fig:s2schema}
\end{figure}

\subsection{估计与检验}

交付点、网格分位点与带宽覆盖点在同一套楔形交算法上计算，种子仍为2026。第一站取$S_1=(0,0)$、示向度$35^\circ$，在楔内作径向--角度样本，样本容量$n=70$。灵敏度固定轴向距离$r=1100$ m。横向偏移依次取$200$、$350$、$500$、$650$、$800$、$1000$ m。楔半角扫描固定在同一轴向的$t=500$ m，另取$0.5^\circ$、$1^\circ$与$1.5^\circ$，对应剖面$(r,t)=(1100,500)$ m，而不是交付点$(800,700)$ m，也不是网格分位点$(800,690.9)$ m。一次清除定义为对径圆覆盖且$D\le 40$ m，或最小包围圆半径本身不超过20 m。

\subsection{结果}

由表\ref{tab:q2des}可知，网格分位点取$r=800$ m、$|t|=690.9$ m，事后直径中位数48.27 m，90\%分位88.94 m，交角中位数$93.85^\circ$。同一套楔内径向--角度样本$n=70$，其中1000 m可听$65/70$（0.93），一次清除$15/70$（$\eta_1=0.21$）。该分位点不是交付点：远轴$\psi(1500)\approx 44.63^\circ<45^\circ$。由表\ref{tab:q2cmp}可知，交付点$(800,700)$ m在同一$n=70$样本上中位直径48.29 m、可听$65/70$，并覆盖轴上$\psi\ge 45^\circ$；$[45^\circ,135^\circ]$覆盖点$(800,795)$ m把可听降到$55/70$，中位直径48.82 m。过样本形心的垂线中位直径46.43 m更小，但直径90\%分位93.69 m超过90 m，且不覆盖全轴。网格比较带为
\[
r\in[800,1300]\,\mathrm{m},\qquad |t|\in[459.1,845.5]\,\mathrm{m},
\]
面积约$3.71\times 10^{5}$ m$^2$，其中$(800,459)$已使远轴$\psi<45^\circ$，故该带不是覆盖集。表\ref{tab:q2xy}给出分位点与交付点在$\theta=35^\circ$时的两侧坐标：从$S_1$走到交付点约1063 m。如图\ref{fig:heatD}所示，事后直径中位数在小$t$和大$r$同时出现时迅速变坏，红框为比较带，星号为网格分位点。如图\ref{fig:heatpsi}所示，比较带大致落在$90^\circ$附近的暖色区域。这表明取较大横向偏移是为了把交角留在可观测区间，而不是把第二站贴在第一示向度轴上。

\begin{table}[htbp]
\centering
\caption{问题二网格分位点与比较带（$n=70$，不是覆盖集）}
\label{tab:q2des}
\begin{tabular}{lr}
\toprule
量 & 取值 \\
\midrule
轴向距离 / m & 800 \\
网格横向偏移 / m & 690.9 \\
交付横向偏移 / m & 700 \\
分位点直径中位数 / m & 48.27 \\
分位点直径90\%分位 / m & 88.94 \\
分位点交角中位数 / 度 & 93.85 \\
楔内径向--角度样本 $n$ & 70 \\
分位点1000 m可听 & $65/70$（0.93） \\
分位点一次清除 $\eta_1$ & $15/70$（0.21） \\
比较带$r$下限 / m & 800 \\
比较带$r$上限 / m & 1300 \\
比较带$|t|$下限 / m & 459.1 \\
比较带$|t|$上限 / m & 845.5 \\
\bottomrule
\end{tabular}
\end{table}

\begin{table}[htbp]
\centering
\caption{分位点与交付点在示向度$35^\circ$时的两侧坐标}
\label{tab:q2xy}
\begin{tabular}{lrrr}
\toprule
点 & $x$ / m & $y$ / m & 至$S_1$距离 / m \\
\midrule
网格 $+t=690.9$ & 259.03 & 1024.82 & 1057.05 \\
网格 $-t$ & 1051.61 & $-107.10$ & 1057.05 \\
交付 $+t=700$ & 253.82 & 1032.28 & 1063.01 \\
交付 $-t=700$ & 1056.82 & $-114.56$ & 1063.01 \\
\bottomrule
\end{tabular}
\end{table}

\begin{table}[htbp]
\centering
\caption{同一$n=70$样本上第二站候选的比较（启发式不是方案）}
\label{tab:q2cmp}
\begin{tabular}{lrrrrr}
\toprule
点 & $|t|$ / m & $\psi\ge 45^\circ$ & $[45^\circ,135^\circ]$ & 中位直径 / m & 1000 m可听 \\
\midrule
网格分位 & 690.9 & 0 & 0 & 48.27 & $65/70$ \\
交付（$\psi\ge 45^\circ$） & 700 & 1 & 0 & 48.29 & $65/70$ \\
带宽覆盖 & 795 & 1 & 1 & 48.82 & $55/70$ \\
形心垂线 & --- & 0 & 0 & 46.43 & --- \\
\bottomrule
\end{tabular}
\end{table}

\begin{figure}[htbp]
\centering
\includegraphics[width=0.82\textwidth]{figures/q2_median_D_heatmap.png}
\caption{事后直径中位数随$(r,|t|)$的变化。红框为比较带，星号为网格分位点；交付点在$|t|=700$ m。}
\label{fig:heatD}
\end{figure}

\begin{figure}[htbp]
\centering
\includegraphics[width=0.82\textwidth]{figures/q2_psi_heatmap.png}
\caption{事后交角中位数随$(r,|t|)$的变化。}
\label{fig:heatpsi}
\end{figure}

由表\ref{tab:q2sens}可知，固定$r=1100$ m时，$t=200$ m的直径中位数为93.12 m，$t=800$ m为58.03 m，相差35.09 m。决策因此取较大横向偏移而不是贴轴。$t=1000$ m时1000 m可听比例降到0.07，直径中位数回升到60.29 m，说明横向不能无限加大。如图\ref{fig:sens}所示，左幅中位直径随$t$先降后升，右幅可听比例在$t>800$ m后迅速下降；40 m参考线表明，题设$1^\circ$下单次两站交会通常还不能保证一次清除。如图\ref{fig:errsweep}所示，楔半角扫描固定在$(r,t)=(1100,500)$ m：半角从$0.5^\circ$增到$1.5^\circ$，中位直径从30.99 m增到92.65 m，一次清除从1变为0。该0--1翻转属于灵敏度剖面$(1100,500)$ m；网格分位点$(r,t)=(800,690.9)$ m的一次清除比例为$\eta_1=0.21$，交付点是另一组坐标$(800,700)$ m。这表明在$(1100,500)$ m处，题设$1^\circ$恰好卡在一次清除的边界上。

\begin{table}[htbp]
\centering
\caption{固定$r=1100$ m时横向偏移的灵敏度}
\label{tab:q2sens}
\begin{tabular}{rrrrr}
\toprule
$t$ / m & 中位直径 / m & 90\%分位 / m & 交角中位 / 度 & 1000 m可听 \\
\midrule
200 & 93.12 & 166.34 & 149.85 & 0.93 \\
350 & 70.15 & 114.93 & 134.54 & 0.93 \\
500 & 61.87 & 89.43 & 124.62 & 0.86 \\
650 & 58.96 & 78.12 & 118.01 & 0.79 \\
800 & 58.03 & 74.80 & 113.40 & 0.64 \\
1000 & 60.29 & 74.92 & 109.11 & 0.07 \\
\bottomrule
\end{tabular}
\end{table}

\begin{figure}[htbp]
\centering
\includegraphics[width=0.98\textwidth]{figures/q2_sensitivity.png}
\caption{左：横向偏移对中位直径的影响；右：1000 m可听比例随横向偏移下降。}
\label{fig:sens}
\end{figure}

\begin{figure}[htbp]
\centering
\includegraphics[width=0.72\textwidth]{figures/q2_err_sweep.png}
\caption{楔半角对中位直径与一次清除的影响，剖面$(r,t)=(1100,500)$ m。橙实线为中位直径，绿虚线为一次清除，红虚线为清除直径40 m。}
\label{fig:errsweep}
\end{figure}

\subsection{诊断}

有效接收半径未知，不能保证对$\Omega_1$内每一点都听到；交付点用1500 m上界给出$\psi\ge 45^\circ$覆盖，并用1000 m可听比例报告可听代价。第一站很靠场地边缘时，$r=800$ m可能把$S_2$送出盘外，协议允许盘外检测，移动耗时仍按直线距离计算。网格分位点的一次清除比例$\eta_1=0.21$，因而两站交会通常还不能一次清除。问题三在$R_{\mathrm{sec}}\le 20$ m时按圆心证书清除，一站示向度多数延后到覆盖骨架走完之后再沿测得示向度搜索。$t=200$ m把交角中位数推到约$150^\circ$，几何接近共线，直径中位数超过90 m，该点不能作为覆盖点。$|t|=700$ m由式\eqref{eq:psi45}给出，不是把$690.9$ m圆整；近端$\psi(5)\approx 138.64^\circ$是该引理的残差。问题三全向追逐不把该交付点当作默认第二站。

\section{问题三：全向源的侦听网与楔形追逐}

\subsection{机理}

全向源在覆盖角内处处可听，限制只来自距离。场地半径1800 m，最小有效接收半径1000 m，原点到边缘的距离大于1000 m，因此“站在原点扫20个频道”不是完备策略。侦听网的目标是：对圆盘内任意点$G$，9个侦听点中至少有一个$S$使$\|S-G\|\le 1000$。该最近距离写成
\[
\min\bigl(\|G\|,\ \min_{k=0,\ldots,7}\|G-R_k\|\bigr),
\]
其中$R_k$为半径1000 m的环点：第一项是到原点的距离，第二项是到环点的距离。该最近距离的最大值$\rho$按原点--环点Voronoi顶点半径$\tau\approx 541.196$ m与场地边界角平分点$P$分款：$r\le\tau$最近的是原点，$r>\tau$最近的是环点；1000 m是设计听距，不是这条最近站分款。跳过尚未访问的侦听点，必须按连续残留而不是有限核对点：发现站是实际扫描过剩余静默频道的位姿，残留是场地减去这些闭圆盘。听到示向度之后，源已经被限制在一个$2^\circ$楔里，但发现层仍把该点的剩余静默频道扫完，并不沿该射线立刻离开覆盖骨架。轴上清除点$P$满足$\|SP\|=\|SG\|=r$时，弦长$\|PG\|=2r\sin(|\alpha|/2)$；对一切$|\alpha|\le 1^\circ$保证命中，当且仅当$r\le 20/(2\sin 0.5^\circ)$。$r\tan 1^\circ$是轴向距离$r$处垂直截面半宽，只是保守充分条件。9盘保证存在某个网点距离不超过$\rho$，并不保证首次听到的站满足该距离。沿射线的扩展搜索、横向有限振幅与沿示向度15 m都不是这个弦圆。问题二的交付点不是本问默认第二站。时间模型为
\begin{equation}
\Delta T_{\mathrm{measure}}=\frac{\|P_{\mathrm{new}}-P_{\mathrm{old}}\|}{5}+1\cdot\mathbf{1}_{c\neq c_{\mathrm{radio}}}+5,
\label{eq:tmeas}
\end{equation}
\begin{equation}
\Delta T_{\mathrm{clear}}=\frac{\|P_{\mathrm{new}}-P_{\mathrm{old}}\|}{5}+\begin{cases}5 & \text{success}\\ 3 & \text{no\_target\_in\_range.}\end{cases}
\label{eq:tclr}
\end{equation}
\texttt{/clear}不切换测向机频道。换台代价虽小，但在20个频道上反复空扫会把虚拟时间和现实时间一起耗尽，所以频道顺序优先当前测向机频道；覆盖路点上的静默扫描不因一站示向度而中断。接口动作次数另设闸$A_{\max}=900$：该上限是算法自设，原题未规定。覆盖循环、静默扫描、射线归航与追逐在达到该上限时返回；程序运行时间、测试窗口与技术异常同样中断。这些情形本身不是覆盖空集，也不进入检测全称。

\subsection{模型}

侦听网（listening net）取原点加上半径1000 m的8等分点，共9个点。如图\ref{fig:scouts}所示，左幅为该9点网，右幅为问题四在同一9盘之外加外环12点的检测网。其基本形式为
\[
\mathcal{N}_3=\bigl\{(0,0)\bigr\}\cup\bigl\{1000(\cos\tfrac{k\pi}{4},\sin\tfrac{k\pi}{4}):k=0,\ldots,7\bigr\}.
\]
点位为
\begin{equation}
O=(0,0),\qquad
R_k=\bigl(1000\cos\tfrac{k\pi}{4},\,1000\sin\tfrac{k\pi}{4}\bigr),\quad k=0,\ldots,7.
\label{eq:q3net}
\end{equation}
表\ref{tab:q3pts}按角色列出半径与个数。侦听网$\mathcal{N}=\{O\}\cup\{R_0,\ldots,R_7\}$。对任意$G\in\mathbb{D}(0,1800)$，9盘最近距离为
\begin{equation}
\min_{S\in\mathcal{N}}\|G-S\|=\min\Bigl(\|G\|,\ \min_{k=0,\ldots,7}\|G-R_k\|\Bigr).
\label{eq:q3mindist}
\end{equation}
场地边界角平分点到最近环点的余弦距离为
\begin{equation}
\rho=\sqrt{1800^2+1000^2-2\cdot 1800\cdot 1000\cos(\pi/8)}\approx 956.05\,\mathrm{m}.
\label{eq:q3rho}
\end{equation}
$\rho$是9盘最近距离在场地上的最大值。8环在原点的最近距离为$\min_k\|O-R_k\|=1000>\rho$，故8环单独不能覆盖原点；9盘在原点的最近距离为$0\le\rho$，落入式\eqref{eq:q3mindist}的第一项。原点残差不在8环最小里。

\textbf{引理（全向侦听覆盖，9盘）。}假设场地为闭圆盘$\mathbb{D}(0,1800)$，侦听网为原点加半径1000 m的8等分环，设计接收半径$R_{\mathrm{eff}}=1000$ m。则对任意$G\in\mathbb{D}(0,1800)$，
\[
\min_{S\in\mathcal{N}}\|G-S\|\le\rho\approx 956.05\,\mathrm{m}.
\]
从而9盘覆盖整个场地。该界是覆盖模型；均匀圆盘抽样只作核对。

\textbf{证明。}由旋转对称，取$\arg G=\theta\in[0,\pi/4]$，令$\delta=\min(\theta,\pi/4-\theta)\le\pi/8$。记$\gamma=\cos(\pi/8)$，$\tau=1000/(2\gamma)\approx 541.196$ m。$\tau$是原点与环点在角平分线上的Voronoi顶点半径。分两款。

第一款：若$\|G\|=r\le \tau$，用原点，$\|G-O\|=r\le \tau<\rho$。只写$r\le 1000$只能说明不超过1000 m，尚未说明不超过$\rho$；本款把原点覆盖收到$\tau$，才进入$\rho$的界。

第二款：若$\tau\le r\le 1800$，最近环点与$G$的圆心角为$\delta$。余弦定理
\begin{equation}
\|G-R\|^2=r^2+1000^2-2\cdot r\cdot 1000\cos\delta.
\label{eq:q3cos}
\end{equation}
因$\delta\le\pi/8$，有$\cos\delta\ge \gamma$，从而
\[
\|G-R\|^2\le f(r):=r^2+10^6-2000\,\gamma r.
\]
$f$是开口向上的二次函数，因而凸。于是$f(r)\le\max\{f(\tau),f(1800)\}$。直接计算$\gamma\tau=500$，故$f(\tau)=\tau^2$；又$f(1800)=\rho^2$，且$\tau<\rho$，所以$f(r)\le\rho^2$。这就给出所声称的上界$\rho$。极小点$1000\cos\delta\le 1000$足以说明$f$在$[1000,1800]$上递增，从而场地边界角平分点
\[
P=1800\bigl(\cos\tfrac{\pi}{8},\sin\tfrac{\pi}{8}\bigr)
\]
是该区间上的最远点。证毕。

同一最远点公式对6环给出约$1059.39$ m，超过1000 m，8环不是任意环数。如图\ref{fig:ninedisk}所示，点划线圆是原点--环点Voronoi顶点半径$\tau$，虚线圆是环点所在的1000 m设计听距圆，实线圆是场地；$V$为角平分线上的Voronoi顶点，$P$为最远点，红弦为$\rho$。点$G$落在$(\tau,1000]$内，最近侦听点是环点而不是原点。线段$OR_0$是8环在原点的最近距离1000 m。表\ref{tab:q3cases}按分款列出覆盖侦听点。

记场地$\mathbb{A}=\overline{\mathbb{D}}(0,1800)$，发现站$\mathcal{S}$为实际扫描过剩余静默频道的位姿。连续残留
\begin{equation}
\mathrm{Res}_{\mathrm{cont}}(\mathcal{S})=\bigl\{p\in\mathbb{A}:\ \min_{S\in\mathcal{S}}\|p-S\|>1000\bigr\}
\label{eq:rescont}
\end{equation}
是场地减去这些闭圆盘。侦听点$Q$的透镜$K_Q=\overline{\mathbb{D}}(Q,1000)\cap\mathbb{A}$。其基本形式为：丢掉$Q$当且仅当$K_Q$上发现站的覆盖半径不超过1000 m，即$\overline{\mathbb{D}}(Q,1000)\cap\mathrm{Res}_{\mathrm{cont}}=\emptyset$。

\textbf{引理（连续残留剪枝，全向）。}假设源为全向，$R_{\mathrm{eff}}\ge 1000$ m，发现站都扫描过当时仍静默的频道。则：（i）尚未听到的源落在$\mathrm{Res}_{\mathrm{cont}}$内。（ii）若$Q\in\mathbb{A}$且$\overline{\mathbb{D}}(Q,1000)\cap\mathrm{Res}_{\mathrm{cont}}=\emptyset$，则$Q$听不到任何剩余全向源。（iii）把仍与残留相交的9网点加回覆盖目标，因此保持对残留的覆盖。（iv）若残留为空，剩余9网点不再是覆盖义务。

第（ii）条须用场地凸性，因为实际$R_{\mathrm{eff}}$可以大于1000 m。设$Q\in\mathbb{A}$能听到尚未发现的源$G$，记$d=\|QG\|\le R_G$。若$d\le 1000$，则$G\in K_Q$，已被某发现站的1000 m盘覆盖，与尚未发现矛盾。若$d>1000$，在线段$QG$上取$\|QP\|=1000$；凸性得$P\in K_Q$，存在发现站$S$使$\|SP\|\le 1000$，于是$\|SG\|\le d\le R_G$，该站应已听到$G$，矛盾。

\textbf{引理（9网仍打到非空残留）。}若$\mathrm{Res}_{\mathrm{cont}}\ne\emptyset$，则存在$Q\in\mathcal{N}$与之相交。否则残留含于$\mathbb{A}\setminus\bigcup_{Q\in\mathcal{N}}\overline{\mathbb{D}}(Q,1000)$，与9盘覆盖引理矛盾。

\textbf{引理（透镜上覆盖半径的最大值类）。}设$K$为若干闭圆盘之交，$d_{\mathcal{S}}(p)=\min_{s\in\mathcal{S}}\|p-s\|$。则$\rho(\mathcal{S},K)=\max_{p\in K}d_{\mathcal{S}}(p)$在下列之一取得：$\mathcal{S}$在$K$内的Voronoi顶点；Voronoi边与某边界圆$\partial\overline{\mathbb{D}}(c_j,r_j)$的交且该点在$K$内；一站在边界圆上的径向远交点；$K$的顶点。

\textbf{证明。}$d_{\mathcal{S}}$是凸函数的点态最小。Voronoi胞腔内部$d_{\mathcal{S}}(p)=\|p-s\|$无临界点（除非$p=s$，那是最小）。在Voronoi边的内部，距离平方沿边是严格凸二次函数，因而也不能取得局部最大值；内部候选只剩下Voronoi顶点。在$\partial K$上，$\mathrm{cell}(s)\cap\partial K$的连通块是圆弧；当$s\neq c_j$时$\|p-s\|$在圆上的驻点只在直线$s\,c_j$上，远交点才是最大候选。这些弧的端点是Voronoi边与$\partial K$的交，或$K$的顶点。

退化：$s=c_j$时整条圆周距离恒定（$Q=O$且$\mathcal{S}=\{O\}$即出现），“唯一径向远点”无定义，任取圆周上一点即可代表该常值。重复站点先合并；相切给出单交点；多站共圆时外心重合，该点仍落在Voronoi顶点类。对本问的两盘透镜$K_Q$，平面上内部Voronoi顶点需要三站；两站图只有平分线，最大值在$\partial K$上；一站最大值在远径向、常值圆周或透镜顶点。故上述类穷尽最大值。72角边界抽样只能抬高估计，漏算候选会低估最大值，抽样不能修复完备性。剪枝只在已证明上界$U\le 1000$时执行。由$K\subseteq K^+$只能得$\rho(K)\le\rho(K^+)$，$\rho(K^+)\le 1000+\varepsilon$不能推出$\rho(K)\le 1000$。数值不确定则保留。

有限边界样本不是该残留。如图\ref{fig:q3res}所示，场地边界$G_\star=1800\angle 7.5^\circ$满足$\|G_\star-R_0\|\approx 819.023$ m$<1000$ m，且到其余8个侦听点的最近距离约$1176.405$ m，都大于1000 m。覆盖引理的最远点$P=1800\angle 22.5^\circ$同时落入$R_0$与$R_1$的1000 m盘，
\[
\|P-R_0\|=\|P-R_1\|=\rho\approx 956.051\,\mathrm{m}<1000\,\mathrm{m}.
\]
访问$R_1$即可使该有限点离开残留，却盖不住$G_\star$。八个场地边界角平分点$1800\angle(22.5^\circ+k\cdot 45^\circ)$在访问$\mathcal{N}\setminus\{R_0\}$之后全部落入其余网点的1000 m盘，因而有限角平分样本可以已空；此时$G_\star$仍只落在$R_0$的独有透镜内。若因这些有限点已空而跳过$R_0$，则$G_\star$不会被听到。本策略不用有限点集去剪枝、回退或过滤静默扫描。

机器狗从原点、频道1出发。剩余覆盖目标初始化为9网点。\texttt{/enter}后先在原点扫描静默频道。覆盖循环另设动作闸$A_{\max}=900$：该上限是算法自设，原题未规定。每一次覆盖路点静默扫描、射线归航与收尾追逐都在$A\ge A_{\max}$时返回；程序运行时间、测试窗口与技术异常同样中断。覆盖目标走开放旅行商问题（open travelling salesman problem, open TSP）。其基本形式为：从当前位姿出发，访问每个剩余覆盖路点恰好一次，不强制回到起点；点数不超过11时用Held--Karp精确求解。下一站可以是覆盖路点、共享D-最优站、一站示向度的接合点，或两站及以上的服务点。两站服务点进入开放TSP时不检查清除证书。若已选出共享D-最优站$p_{\mathrm{D}}$，记开放TSP在当前位姿$P_{\mathrm{now}}$下的首站为$p_1$，路程差
\begin{equation}
\Delta_{\mathrm{first}}=\|p_{\mathrm{D}}-P_{\mathrm{now}}\|-\|p_1-P_{\mathrm{now}}\|.
\label{eq:dfirst}
\end{equation}
当$\Delta_{\mathrm{first}}\le 80$ m时，下一站取$p_{\mathrm{D}}$。当$p_{\mathrm{D}}$落在剩余覆盖路点的120 m邻域内或覆盖某个剩余透镜，并且$\Delta_{\mathrm{first}}\le 220$ m时，下一站同样取$p_{\mathrm{D}}$而不取$p_1$。该抢先规则只改变访问顺序，不删除残留非空的9网透镜。最近剩余、D-最优与边上顺路探测是时间启发式：它们只在路点同样连续覆盖该透镜时替换侦听点，证书回退仍是把与残留相交的9网点加回。覆盖路点对尚未清除的静默频道做检测时不进入追逐：\texttt{near}就地清除，\texttt{direction}只记示向度，该点把剩余静默频道扫完后再走下一覆盖路点。题目约束$N\in[10,16]$：已听到16个不同频道后，剩余静默频道不可能再是源，这是假设而不是覆盖引理。

\textbf{引理（圆心再测，全向）。}假设源为全向，$R_{\mathrm{eff}}\ge 1000$ m，且$R_{\mathrm{sec}}(\Omega)\le 1000$ m。则$G\in\Omega\subseteq\overline{\mathbb{D}}(C,R_{\mathrm{sec}})$，在最小包围圆圆心测量不是\texttt{no\_signal}。该引理对定向瓣不成立。圆心\texttt{no\_signal}后的18 m六邻点只在$R_{\mathrm{sec}}\le 80$ m时尝试，并且不是半径80 m盘的覆盖。

\textbf{引理（轴上弦长清除）。}假设$G$落在以$S$为顶点、半角$1^\circ$的闭楔内，实际角误差为$\alpha$，记$r=\|G-S\|$。清除点$P$位于测得示向度轴上，且$\|SP\|=\|SG\|=r$。则弦长
\begin{equation}
\|PG\|=2r\sin\bigl(|\alpha|/2\bigr).
\label{eq:chord}
\end{equation}
对某个实际源，命中当且仅当$2r\sin(|\alpha|/2)\le 20$。对一切允许误差$|\alpha|\le 1^\circ$，保证命中当且仅当
\begin{equation}
2r\sin 0.5^\circ\le 20\qquad\Leftrightarrow\qquad r\le r_{\mathrm{ray}}=\frac{20}{2\sin 0.5^\circ}\approx 1145.930\,\mathrm{m}.
\label{eq:rray}
\end{equation}
$r\tan 1^\circ$是轴向距离为$r$的垂直截面半宽，与弦长不同。它可以给出保守充分条件，不能写成上述充要；其数值$20/\tan 1^\circ\approx 1145.799$ m。

9盘只保证存在$S\in\mathcal{N}$使$\|G-S\|\le\rho$。这不是首次听到的站。反例$G=(1400,0)$、$R_{\mathrm{eff}}=1500$ m：原点首次听到，距离1400 m，$1^\circ$误差下$2800\sin 0.5^\circ\approx 24.43>20$。存在一个原始网点，在真实距离已知的理想条件下该轴上清除圆能盖住源；真实距离未知，这仍不是可执行的清除保证。可执行的清除保证仍是$R_{\mathrm{sec}}\le 20$ m（或对径圆覆盖且$D/2\le 20$ m）的圆心证书。该命题的覆盖对象只是这个弦圆。沿射线按步长$70$ m、$120$ m作扩展搜索、越过源后的一维收缩、横向有限振幅以及沿示向度再\texttt{/clear} 15 m，都不是这个圆。距离由覆盖结束后的单向扩展搜索得到：前向听到则继续前进，\texttt{near}则就地清除，\texttt{no\_signal}或示向度翻转则认为已越过源，在最后前向站与越源站之间作一维收缩后再清除。楔在$r=1500$ m处的垂直截面半宽为
\begin{equation}
R_{\mathrm{eff}}^{\max}\tan 1^\circ=1500\tan 1^\circ\approx 26.18\,\mathrm{m}.
\label{eq:sliver}
\end{equation}
横向测点取离散表
\[
\{\pm 12,\pm 18,\pm 26,\pm 34\}\,\mathrm{m}.
\]
其中$\pm 34$ m大于式\eqref{eq:sliver}的26.18 m，落在1500 m楔截面之外；26.18 m是该截面半宽，不是这张表的上界。该表是越过源之后的一维搜索启发式，不是一切$r$上的连续覆盖。问题二的交付点$(800,700)$ m不是本问默认第二站。

覆盖骨架尚未走完时，一站示向度不加入当前路径，除非接合点距当前位姿不超过420 m，或第一测向站到原点的距离大于200 m且该接合点能连续割完某个剩余9网透镜。半径200 m是该并入规则的阈值：在原点听到，以及在原点200 m闭圆盘内听到，都不满足该条。两站及以上的服务点不受该阈值约束：只要该频道已有不少于两站示向度，服务点就进入当前开放TSP。

\textbf{策略（一站示向度并入覆盖路径）。}假设覆盖骨架尚未走完，该频道仅有一站示向度，接合点$p$为后验形心或沿测得示向度550 m处，第一测向站为$S^{(1)}$。则$p$进入当前开放TSP，当且仅当下列之一成立：（i）$\|p-P_{\mathrm{now}}\|\le 420$ m；（ii）$\|S^{(1)}\|>200$ m，且$p$的1000 m盘能连续割完某个剩余9网透镜。覆盖已走完时，该频道改走服务TSP。该条是时间启发式，不是检测定理。

\textbf{策略（两站服务点并入覆盖路径）。}假设该频道已记录\texttt{near}，或已有不少于两站示向度。记服务点$p$为最小包围圆圆心；若定位多边形不足三点，则取后验形心或沿最新示向度550 m处。则$p$进入当前开放TSP，不论$R_{\mathrm{sec}}$是否已不超过20 m。选中该城市后，先按圆心再测引理处理定位多边形：证书半径则在该圆心\texttt{/clear}；若得到新示向度则下一轮重算服务点；若返回开放，再沿最新示向度作单向归航。归航入口若已记录\texttt{near}，则就地清除。覆盖骨架尚未走完时，该条同样执行。清除证书仍是$R_{\mathrm{sec}}\le 20$ m（或对径圆覆盖且$D/2\le 20$ m）；未达证书并不阻止该服务点进入开放TSP。近邻已达证书的圆心可另作顺路清除，那是另一入口，不是本条的准入条件。该条是时间启发式，不是检测定理。

边上顺路探测采用Clarke--Wright节省法则（Clarke--Wright savings）\cite{clarke1964}。其基本形式为把点$p$插入边$ab$的额外路程
\begin{equation}
\Delta_{\mathrm{edge}}(p;a,b)=\|p-a\|+\|p-b\|-\|a-b\|.
\label{eq:cwextra}
\end{equation}
仅当$\Delta_{\mathrm{edge}}\le 56$ m且示向度与该覆盖边的夹角不超过$8^\circ$时，才在该边上测一站，并不把该频道改写成覆盖城市。\texttt{near}以及两站及以上的服务点进入当前开放TSP，不要求该频道已达清除证书；选中后先圆心再测，仍开放则沿测得示向度归航。覆盖骨架尚未走完时，这条两站路径同样执行。时间启发式失败时，覆盖义务仍是与连续残留相交的9网点，按开放TSP从当前位姿重排；覆盖骨架本身不改走命中即射线的星形。两站归航是服务层动作，不是把覆盖义务改写成星形。

若干一站示向度可以共用一个D-最优（D-optimal）第二站\cite{bishop2010}。两站方位的Fisher信息矩阵（Fisher information matrix, FIM）满足$\det(F_S+F_P)\propto\sin^2\psi/(r_S^2 r_P^2)$。其基本形式为：对第一示向度沿轴的先验距离$\mathcal{R}=\{250,450,650,850,1100\}$ m取期望对数行列式
\begin{equation}
I(p)=\sum_{c}\frac{1}{|\mathcal{R}|}\sum_{r\in\mathcal{R}}\bigl(2\log\sin\psi_c(r)-2\log r_S-2\log r_P\bigr),
\label{eq:dopt}
\end{equation}
并统计期望交会已把保守半径$r_S\tan 1^\circ/\sin\psi$压到28 m以内的频道数。同站示向度按圆心角间隙大于$45^\circ$分簇，扇形宽度超过$75^\circ$再切开。把候选站$p$插入当前覆盖骨架（当前位姿出发的开放路径）的最便宜插入额外路程记为$\Delta_{\mathrm{ins}}(p)$：在骨架各边上按式\eqref{eq:cwextra}计算插入增量，并与接到骨架终点的路程比较，取较小者。

\textbf{策略（共享D-最优专用切割）。}假设候选站$p$用于压缩不少于两个一站楔。若$p$已落在某个剩余覆盖路点的120 m邻域内，或$p$本身覆盖某个剩余透镜，则不再用额外路程闸，只要求期望压到28 m的频道数不少于2，或至少压到1个频道且可听频道数不少于2。否则须先有不少于2个频道被期望压到28 m以内。覆盖骨架已空时，再要求$\|p-P_{\mathrm{now}}\|\le 920$ m。覆盖骨架非空时，接受该专用切割当且仅当下列之一成立：（i）$\Delta_{\mathrm{ins}}(p)\le 120$ m且$\|p-P_{\mathrm{now}}\|\le 920$ m；（ii）期望交会已把不少于3个频道压到28 m以内，$\Delta_{\mathrm{ins}}(p)\le 280$ m，且$\|p-P_{\mathrm{now}}\|\le 920$ m；（iii）$p$连续割完某个剩余透镜且$\Delta_{\mathrm{ins}}(p)\le 520$ m。第（ii）款是三射线扇的一站内切割，额外280 m仍短于一次环间跳跃。该条是时间启发式：共享站上只测向，不从该站对每个源作星形清除。两站服务点进入覆盖TSP时不要求已达证书；选中后的圆心再测与射线归航属于服务层，覆盖义务仍是与残留相交的9网点。

\textbf{定理（全向检测，正常结束）。}假设问题三：源为全向，$G\in\mathbb{A}$，$R_{\mathrm{eff}}\ge 1000$ m，$N\in[10,16]$，并且覆盖循环按前三类退出结束（残留空、静默频道空、或已听满16个）。则每个全向源都已被听到。这是正常结束时的检测正确性，不是题面“确保全部清除”，也不证明覆盖循环必在有限步或$A<900$内走完。

证明：反证。若某源$G$从未被听到且循环仍按前三类结束，则每一次静默频道检测（因而每一个发现站）满足$\|G-S\|>R_{\mathrm{eff}}\ge 1000$，故该源始终属于$\mathrm{Res}_{\mathrm{cont}}$。9盘覆盖把它放进某个透镜$K_Q$。连续残留引理与透镜上覆盖半径最大值类保持该$Q$（或其连续覆盖路点）为覆盖目标，直到一次静默扫描盖住$K_Q$，从而听到该源。前三类退出都与“仍有未听到的源”矛盾。证毕。

覆盖循环还有其它终止，见表\ref{tab:q3exit}：$A_{\max}=900$是算法自设闸，原题未规定该上限；程序运行时间、测试窗口与技术异常同样中断。这些残差下本定理不声称已听到每一个源。清除证书仍是$R_{\mathrm{sec}}\le 20$ m（或对径圆覆盖且$D/2\le 20$ m），圆心再测引理的前提是$G\in\Omega$。本文不补“加密测站直到半径不超过20 m必在有限步收敛”的证明。同构几何核验24个场景的动作次数介于123与213之间，演练三次为132、150、160，都低于900，只说明这些样本没有碰到动作闸，不能把该闸从残差里删掉，也不是本定理。覆盖骨架的访问顺序、D-最优插入与服务TSP都不删除残留非空的9网透镜，因此也不进入本定理。

\begin{table}[htbp]
\centering
\caption{问题三覆盖循环的退出集}
\label{tab:q3exit}
\begin{tabular}{>{\raggedright\arraybackslash}p{0.42\textwidth}>{\raggedright\arraybackslash}p{0.50\textwidth}}
\toprule
退出集 & 与尚未听到的源的关系 \\
\midrule
$\mathrm{Res}_{\mathrm{cont}}=\emptyset$ & 与该源始终属于连续残留矛盾 \\
剩余静默频道为空 & 该源已有示向度，已被听到 \\
已听到16个不同频道 & 由$N\le 16$，不存在尚未听到的源 \\
$A\ge A_{\max}=900$ & 算法自设闸，原题未规定；覆盖扫描可以尚未完成 \\
程序运行时间、测试窗口、技术异常 & 额外残差：定理不声称已听到每一个源 \\
\bottomrule
\end{tabular}
\end{table}

覆盖骨架走完之后，尚未清除的频道走一条服务TSP：接合点取后验形心或沿示向度550 m处；随后对仍开放的频道沿最新示向度作单向搜索。覆盖骨架尚未走完时，两站服务点一旦被开放TSP选中，同样先圆心再测，仍开放则沿测得示向度归航，并不等到骨架走完。全向局部化见表\ref{tab:q3chase}，时间层插入条件见表\ref{tab:q3insert}。两站及以上先按圆心再测引理处理定位多边形：若$R_{\mathrm{sec}}\le 20$ m，或对径圆覆盖且$D/2\le 20$ m，在该圆心\texttt{/clear}；若$20<R_{\mathrm{sec}}\le 1000$ m，在最小包围圆圆心再测，\texttt{near}则清除，\texttt{direction}则新楔并入$\Omega$并重算，\texttt{no\_signal}仅当$R_{\mathrm{sec}}\le 80$ m时尝试圆心\texttt{/clear}与18 m六邻点（不覆盖80 m盘）。圆心再测后仍开放，才沿最新示向度作单向搜索。沿射线归航结束后若仍未清除且动作预算尚未用尽，在最后命中点沿示向度再\texttt{/clear} 15 m。该步并不等到动作预算用尽才执行，也不是覆盖证书。

如图\ref{fig:coverfirst}所示，左幅覆盖骨架沿9网点走开放TSP，右幅把一站示向度延后到覆盖结束之后，接合射线离开覆盖辐条。如图\ref{fig:chase}所示，左幅为一站楔，中幅为服务层沿测得示向度的单向测站，右幅为轴上弦长$\|PG\|=2r\sin(|\alpha|/2)$与半径20 m的清除圆。左幅楔半角为示意放大，正文始终取$1^\circ$。每步使用新的请求编号，串行等待。

\begin{table}[htbp]
\centering
\caption{问题三全向局部化：证书、圆心再测与射线搜索}
\label{tab:q3chase}
\begin{tabular}{>{\raggedright\arraybackslash}p{0.46\textwidth}>{\raggedright\arraybackslash}p{0.46\textwidth}}
\toprule
条件 & 动作 \\
\midrule
覆盖路点上的\texttt{near} & 就地清除 \\
覆盖路点上的\texttt{direction} & 只记示向度，扫完该点静默频道 \\
$R_{\mathrm{sec}}\le 20$ m，或对径圆覆盖且$D/2\le 20$ m & 在该圆心\texttt{/clear}（清除证书） \\
$20<R_{\mathrm{sec}}\le 1000$ m & 在最小包围圆圆心再测 \\
圆心\texttt{no\_signal}且$R_{\mathrm{sec}}\le 80$ m & 圆心\texttt{/clear}与18 m六邻点（不覆盖80 m盘） \\
圆心再测后仍开放，或仅一站且已进入服务层 & 沿最新示向度单向扩展搜索 \\
覆盖未完且两站服务点被选中 & 先圆心再测，仍开放则沿测得示向度归航 \\
\texttt{no\_signal}或示向度翻转 & 在最后前向站与越源站之间一维收缩，再\texttt{/clear} \\
横向有限振幅 & 离散表$\{\pm 12,\pm 18,\pm 26,\pm 34\}$ m；26.18 m是1500 m截面半宽，不是该表上界 \\
射线归航后仍未清除、预算未尽 & 最后命中点沿示向度再\texttt{/clear} 15 m（非证书） \\
\bottomrule
\end{tabular}
\end{table}

\begin{table}[htbp]
\centering
\caption{问题三时间层：覆盖骨架上的插入条件}
\label{tab:q3insert}
\begin{tabular}{>{\raggedright\arraybackslash}p{0.46\textwidth}>{\raggedright\arraybackslash}p{0.46\textwidth}}
\toprule
对象 & 进入当前路径的条件 \\
\midrule
覆盖路点 & 透镜与连续残留相交；开放TSP从当前位姿重排 \\
\texttt{near}或两站及以上服务点（覆盖未完） & 服务点进入当前开放TSP，不要求已达清除证书；选中后先圆心再测，仍开放则沿测得示向度归航 \\
已达证书的近邻圆心 & 顺路清除；不是两站入口的准入条件 \\
一站示向度（覆盖未完） & 接合点距当前$\le 420$ m，或第一测向站$\|S^{(1)}\|>200$ m且接合点割完剩余透镜 \\
D-最优下一站抢先 & $\Delta_{\mathrm{first}}\le 80$ m；或该站落在覆盖路点120 m邻域或覆盖剩余透镜，且$\Delta_{\mathrm{first}}\le 220$ m \\
边上顺路探测 & $\Delta_{\mathrm{edge}}\le 56$ m且示向度与该边夹角$\le 8^\circ$；不改写成覆盖城市 \\
共享D-最优专用切割 & 额外路程$\le 120$ m且到当前$\le 920$ m；三频道压到28 m时额外路程$\le 280$ m；连续割完剩余透镜时额外路程$\le 520$ m；已在覆盖路点120 m邻域或覆盖剩余透镜则不再用额外路程闸 \\
覆盖未完且启发式失败 & 把仍与残留相交的9网点加回，按开放TSP走 \\
覆盖循环动作闸 & $A_{\max}=900$为算法自设；达到上限、超时或异常则返回，残留透镜可以尚未扫描 \\
覆盖已完的一站示向度 & 后验形心或沿示向度550 m处，走一条服务TSP \\
\bottomrule
\end{tabular}
\end{table}

\begin{table}[htbp]
\centering
\caption{连续残留上覆盖半径的最大值类}
\label{tab:q3rho}
\begin{tabular}{ll}
\toprule
类 & 几何对象 \\
\midrule
Voronoi顶点 & 三站外心，空外接圆 \\
Voronoi边与$\partial K$ & 两站垂直平分线交边界圆 \\
径向远极值 & 一站与圆心连线的远交点 \\
透镜顶点 & 两边界圆的交点 \\
Voronoi边内部 & 距离平方沿边严格凸，不能取局部最大 \\
同心圆常值 & $s=c_j$时整圆距离恒定，$Q=O$且$\mathcal{S}=\{O\}$即出现 \\
重复站点、相切、共圆 & 先合并；单交点；外心重合仍属顶点类 \\
72角抽样 & 只能抬高估计，不是完备性 \\
\bottomrule
\end{tabular}
\end{table}

\begin{figure}[htbp]
\centering
\includegraphics[width=0.62\textwidth]{figures/q3_nine_disk.png}
\caption{问题三9盘覆盖。实线为场地1800 m，虚线为环点所在的1000 m设计听距圆，点划线为原点--环点Voronoi顶点半径$\tau\approx 541$ m。$r\le\tau$由原点覆盖；$r>\tau$由最近环点覆盖，其中$r\in(\tau,1000]$的最近侦听点是环点而不是原点。角平分线上的方点$V$为Voronoi顶点，最远点$P$在场地边界角平分处，红弦为$\rho$。白点$G$落在$(\tau,1000]$内。线段$OR_0$长1000 m，故8环不能覆盖原点。}
\label{fig:ninedisk}
\end{figure}

\begin{figure}[htbp]
\centering
\includegraphics[width=0.62\textwidth]{figures/q3_residual.png}
\caption{连续残留与有限样本的差别。红实线圆为$R_0$的1000 m盘，红弦连到只被$R_0$听到的$G_\star=1800\angle 7.5^\circ$；灰虚线圆为$R_1$的1000 m盘，灰弦连到角平分点$P=1800\angle 22.5^\circ$。$P$可被$R_1$盖住，而$G_\star$落在该虚线盘外，仍属$R_0$的独有透镜。有限角平分点已空不能代替$\mathrm{Res}_{\mathrm{cont}}$，否则会错误跳过$R_0$。}
\label{fig:q3res}
\end{figure}

\begin{table}[htbp]
\centering
\caption{有限角平分样本与$R_0$独有透镜}
\label{tab:q3respts}
\begin{tabular}{lcccc}
\toprule
点 & 方位 & 到$R_0$/m & 到$R_1$/m & 到原点/m \\
\midrule
$G_\star$ & $7.5^\circ$ & 819.023 & 1176.405 & 1800 \\
$P$ & $22.5^\circ$ & 956.051 & 956.051 & 1800 \\
\bottomrule
\end{tabular}
\end{table}

\begin{table}[htbp]
\centering
\caption{问题三9盘覆盖的分款}
\label{tab:q3cases}
\begin{tabular}{lll}
\toprule
位置 & 覆盖侦听点 & 最近距离 \\
\midrule
$r\le \tau\approx 541.196$ m & 原点 & $r\le \tau<\rho$ \\
$\tau<r\le 1000$ m & 最近环点（不是原点） & $\|G-R\|^2\le f(r)\le\rho^2$ \\
$\tau\le r\le 1800$ m & 最近环点 & $\|G-R\|^2\le f(r)\le\rho^2$ \\
边界角平分点 $P$ & 两相邻环点 & $\rho\approx 956.05$ m \\
仅写$r\le 1000$ m & 原点 & 不超过1000，尚未说明不超过$\rho$ \\
原点上的8环 & 最近环点（不含原点） & $1000>\rho$ \\
原点上的9盘 & 原点 & $0\le\rho$ \\
\bottomrule
\end{tabular}
\end{table}

表\ref{tab:q3xy}按式\eqref{eq:q3net}列出9个侦听点的坐标，供核验点位。覆盖目标走开放TSP，该表不是固定环序策略。

\begin{table}[htbp]
\centering
\caption{问题三侦听网的点位构成}
\label{tab:q3pts}
\begin{tabular}{lccc}
\toprule
角色 & 半径 / m & 个数 & 方位 \\
\midrule
原点 & 0 & 1 & --- \\
侦听环 & 1000 & 8 & $k\pi/4$ \\
合计 & --- & 9 & --- \\
\bottomrule
\end{tabular}
\end{table}

\begin{table}[htbp]
\centering
\caption{问题三侦听点坐标}
\label{tab:q3xy}
\begin{tabular}{crrl}
\toprule
编号 & $x$ / m & $y$ / m & 角色 \\
\midrule
1 & 0.0 & 0.0 & 原点 \\
2 & 1000.0 & 0.0 & 环 $k=0$ \\
3 & 707.1 & 707.1 & 环 $k=1$ \\
4 & 0.0 & 1000.0 & 环 $k=2$ \\
5 & $-707.1$ & 707.1 & 环 $k=3$ \\
6 & $-1000.0$ & 0.0 & 环 $k=4$ \\
7 & $-707.1$ & $-707.1$ & 环 $k=5$ \\
8 & 0.0 & $-1000.0$ & 环 $k=6$ \\
9 & 707.1 & $-707.1$ & 环 $k=7$ \\
\bottomrule
\end{tabular}
\end{table}

\begin{figure}[htbp]
\centering
\includegraphics[width=0.98\textwidth]{figures/q34_scout_nets.png}
\caption{左：问题三的9点侦听网（点在1000 m圆上）；右：问题四检测网$N_9\cup R_{12}$，外环12个方点在半径2000 m的圆上。}
\label{fig:scouts}
\end{figure}

\begin{figure}[htbp]
\centering
\includegraphics[width=0.98\textwidth]{figures/covering_skeleton.png}
\caption{问题三时间层。左幅覆盖骨架是原点加8环的开放TSP；启发式失败时覆盖义务仍走这9网点，不把覆盖骨架改写成命中即射线的星形。共享D-最优站可按相对TSP首站的路程差抢先作为下一站。右幅一站示向度多数延后到覆盖结束后，接合点取沿测得示向度550 m处。两站服务点在覆盖未完时也可进入同一开放TSP，选中后走圆心再测与射线归航。}
\label{fig:coverfirst}
\end{figure}

\begin{figure}[htbp]
\centering
\includegraphics[width=0.98\textwidth]{figures/chase_flow.png}
\caption{问题三全向局部化几何。左幅楔半角为示意放大，正文取$1^\circ$；中幅为覆盖结束后沿测得示向度的单向扩展搜索；右幅轴上$P$满足$\|SP\|=\|SG\|=r$，弦长$2r\sin(|\alpha|/2)$，清除圆半径20 m。$r\tan 1^\circ$是垂直截面半宽，不是该充要。该图不是发现层的覆盖骨架，也不表示首次听到的站。}
\label{fig:chase}
\end{figure}

\subsection{估计与检验}

覆盖模型是式\eqref{eq:q3mindist}--\eqref{eq:q3rho}的9盘闭式界加上式\eqref{eq:rescont}的连续残留剪枝：$\rho\approx 956.05$ m$<1000$ m。分款以$\tau=1000/(2\cos(\pi/8))\approx 541.196$ m为界：$r\le \tau$用原点，$\tau\le r\le 1800$用最近环点并借$f$的凸性得到$\|G-R\|^2\le\rho^2$。只写$r\le 1000$不能单独推出不超过$\rho$。8环在原点的最近距离为1000 m，大于$\rho$，故不能把$\rho$写成8环最小。场地圆盘上的均匀样本4000个、种子2026，只对该界作抽样核对，核对覆盖率为1。连续残留的最大值类见表\ref{tab:q3rho}；表\ref{tab:q3respts}中的角平分点$P$可被$R_1$盖住，不是连续残留，也不能用来丢掉$R_0$。轴上弦长清除由式\eqref{eq:chord}--\eqref{eq:rray}给出，覆盖对象是$\|SP\|=\|SG\|=r$时的弦圆。时间层插入条件见表\ref{tab:q3insert}：覆盖骨架开放TSP、式\eqref{eq:dfirst}的D-最优抢先（路程差不超过80 m，或该站覆盖剩余透镜且路程差不超过220 m）、三射线专用切割额外路程280 m、一站示向度的200 m第一测向站阈值、两站服务点不论证书进入开放TSP并在选中后圆心再测再归航，以及覆盖未完时把9网点加回。覆盖骨架未走完时两站归航是服务层动作，覆盖义务回退仍是9网点。覆盖循环的退出集见表\ref{tab:q3exit}：前三类给出检测正确性；$A_{\max}=900$是算法自设闸，程序运行时间、测试窗口与技术异常是额外残差，不能用样本动作次数把它们删掉。同构几何核验24个场景全部清除、均值272.227 s/源、路程12.4 km，以及演练全部清除，都只是展示，不是覆盖定理。演练在模拟器窗口连续取三个独立案例：每次结束时完成面板返回干扰源总数，以该总数为分母计算清除比例。正式测试见表\ref{tab:official3}，按题目表1的五列填写案例编码、HTTP清除个数、平均定位清除时间与程序运行时间；完成弹窗不给出总数，不得把该清除个数写成干扰源总数。不把演练行填入该表。

\subsection{结果}

由表\ref{tab:q3cov}可知，9个侦听点的覆盖半径$\rho\approx 956.05$ m，小于设计接收半径1000 m；4000个均匀样本上的抽样核对覆盖率为1。由表\ref{tab:q3cases}与表\ref{tab:q3cov}可知，8环在原点的最近距离为1000 m，大于$\rho$，9盘在原点的最近距离为0；原点--环点Voronoi顶点半径$\tau\approx 541.196$ m小于$\rho$。如图\ref{fig:ninedisk}所示，最远点$P$由环点覆盖，$r\le\tau$由原点覆盖，$r\in(\tau,1000]$的最近侦听点是环点而不是原点；1000 m虚线是环点所在的设计听距圆，不是$\rho$的分款。由表\ref{tab:q3respts}与图\ref{fig:q3res}可知，$P$到$R_0$、$R_1$的距离都等于$\rho\approx 956.051$ m，而$G_\star$到$R_1$约1176.405 m，超过1000 m：访问$R_1$可使该角平分样本离开残留，却盖不住$G_\star$，故有限角平分点已空不能代替连续残留。如图\ref{fig:covmap}所示，左幅把覆盖样本画成绿色、未覆盖样本画成红色，在该网下未覆盖点不出现。由表\ref{tab:q3cov}与式\eqref{eq:rray}可知，对一切$|\alpha|\le 1^\circ$的轴上弦长清除半径约为1145.930 m；$r\tan 1^\circ$给出的垂直截面充分界约为1145.799 m。9盘只保证存在某个网点距离不超过$\rho$，并不保证首次听到的站落在该弦长界内。如图\ref{fig:coverfirst}所示，覆盖骨架与延后服务射线不是同一条路径。横向有限振幅、开放TSP、D-最优插入及其抢先规则是时间启发式，不能单独证明缩短了移动时间。

\begin{table}[htbp]
\centering
\caption{问题三侦听覆盖}
\label{tab:q3cov}
\begin{tabular}{lr}
\toprule
量 & 取值 \\
\midrule
侦听点个数 & 9 \\
覆盖半径 $\rho$ / m & 956.05 \\
4000点抽样核对 & 1 \\
8环在原点的最近距离 / m & 1000 \\
8环原点是否$\le\rho$ & 0 \\
9盘在原点的最近距离 / m & 0 \\
9盘原点是否$\le\rho$ & 1 \\
场地半径 / m & 1800 \\
设计接收半径 / m & 1000 \\
原点--环点Voronoi半径 $\tau$ / m & 541.196 \\
轴上弦长清除半径 $r_{\mathrm{ray}}$ / m & 1145.930 \\
垂直截面充分界 $20/\tan 1^\circ$ / m & 1145.799 \\
\bottomrule
\end{tabular}
\end{table}

\begin{figure}[htbp]
\centering
\includegraphics[width=0.98\textwidth]{figures/q34_coverage_scatter.png}
\caption{左：问题三9点网的全向覆盖（场地圆盘均匀样本$n=4000$，未覆盖点不出现）；右：问题四检测网$N_9\cup R_{12}$上800个位置--朝向样本，绿色为可听、红色为开背瓣残留。该图是抽样核对，不是$\forall u$覆盖定理。}
\label{fig:covmap}
\end{figure}

由表\ref{tab:q3reh}可知，三次演练的干扰源总数分别为16、12、11，均落在题目给定的$[10,16]$内，且全部为全向源。三次均把全部源清除，比例为1。动作次数为132、150、160，平均定位清除时间为197.42 s、280.02 s、295.00 s，路程为12.0 km、12.5 km、11.6 km。合并不配对均值为250.36 s/源。如图\ref{fig:q3reh}所示，真实个数与清除个数三次重合。这表明在这三个独立案例上，9点网加连续残留听得到，覆盖骨架走完后再按圆心证书或射线搜索收到清除圆内。虚拟时间分别为3158.7 s、3360.2 s、3245.0 s，远长于20分钟的现实程序时限，这与“虚拟时钟由移动与检测累加、现实时钟只约束接口存活”的协议一致。同构几何核验24个场景全部清除，均值272.227 s/源、路程12.4 km，动作次数介于123与213之间，见表\ref{tab:q3twin}。这些样本的动作次数都低于上限900，只说明它们没有碰到算法自设闸；该核验与演练全部清除都不是覆盖定理。

\begin{table}[htbp]
\centering
\caption{问题三演练测试统计（连续残留策略；不是正式测试）}
\label{tab:q3reh}
\scriptsize
\begin{tabular}{llrrrrr}
\toprule
演练 & 案例编码 & 总数 & 清除 & 比例 & 动作 & 平均时间 / s \\
\midrule
1 & UUDJ-WJGP-NBRH-8V2N & 16 & 16 & 1.000 & 132 & 197.42 \\
2 & XUXT-9646-NF9K-42J4 & 12 & 12 & 1.000 & 150 & 280.02 \\
3 & 2GA6-9ADV-XD3E-AVMK & 11 & 11 & 1.000 & 160 & 295.00 \\
\bottomrule
\end{tabular}
\end{table}

\begin{table}[htbp]
\centering
\caption{问题三演练的虚拟时间与路程}
\label{tab:q3virt}
\begin{tabular}{lrrr}
\toprule
演练 & 虚拟时间 / s & 平均定位清除时间 / s & 路程 / m \\
\midrule
1 & 3158.714 & 197.42 & 11993.6 \\
2 & 3360.203 & 280.02 & 12456.0 \\
3 & 3245.032 & 295.00 & 11575.2 \\
\bottomrule
\end{tabular}
\end{table}

\begin{table}[htbp]
\centering
\caption{同构几何核验（非正式测试，不是覆盖定理）}
\label{tab:q3twin}
\begin{tabular}{lr}
\toprule
量 & 取值 \\
\midrule
场景数 & 24 \\
全部清除数 & 24 \\
均值 / (s/源) & 272.227 \\
路程均值 / m & 12442 \\
动作次数最小 & 123 \\
动作次数最大 & 213 \\
动作预算上限 $A_{\max}$ & 900 \\
\bottomrule
\end{tabular}
\end{table}

\begin{figure}[htbp]
\centering
\includegraphics[width=0.98\textwidth]{figures/q3_rehearsal_bars.png}
\caption{问题三三次演练：真实个数与清除个数重合；右幅为平均定位清除时间。横轴为演练1--3。}
\label{fig:q3reh}
\end{figure}

题目要求按表1报告三次正式测试。表\ref{tab:official3}按模拟器完成弹窗的案例编码填写；清除个数为HTTP成功\texttt{/clear}次数，平均定位清除时间为虚拟时间与该次数之比，程序运行时间为\texttt{/enter}与\texttt{/exit}的\texttt{real\_timestamp\_ms}之差。正式弹窗不给出干扰源总数，不得把该清除个数写成总数。不把表\ref{tab:q3reh}的演练行抄入该表。

\begin{table}[htbp]
\centering
\caption{问题三正式测试结果（按题目表1格式；完成弹窗未给出总数）}
\label{tab:official3}
\begin{tabular}{lcccc}
\toprule
测试 & 案例编码 & 清除干扰源个数 & 平均定位清除时间 / s & 程序运行时间 / s \\
\midrule
测试1 & TKUB-AWDW-FWUJ-53UT & 12 & 252.466 & 4.890 \\
测试2 & YZRC-HM6Q-TAFH-W5W9 & 11 & 279.556 & 5.449 \\
测试3 & YYM7-3QP9-TJEH-VBD9 & 13 & 255.325 & 4.136 \\
\bottomrule
\end{tabular}
\end{table}

由表\ref{tab:official3}可知，三次正式测试的案例编码分别为TKUB-AWDW-FWUJ-53UT、YZRC-HM6Q-TAFH-W5W9、YYM7-3QP9-TJEH-VBD9。HTTP成功\texttt{/clear}次数分别为12、11、13，平均定位清除时间分别为252.466 s、279.556 s、255.325 s，程序运行时间分别为4.890 s、5.449 s、4.136 s。完成弹窗不给出干扰源总数，因而这三行不得读成12、11、13个源的全清比例，也不得与表\ref{tab:q3reh}混填。三次合计HTTP清除36次、虚拟时间9423.932 s，合并不配对261.776 s/已清。程序运行时间由\texttt{/enter}与\texttt{/exit}的\texttt{real\_timestamp\_ms}之差给出。命名实现为同一追逐程序、问题三关闭混合源：覆盖路点扫完静默频道后再走，\texttt{near}就地清除，\texttt{direction}不立刻沿射线；时间层为开放TSP、D-最优抢先80/220 m、三射线切割280 m、一站离原点大于200 m才并入，两站未完覆盖也进入开放TSP后圆心再测。

\FloatBarrier
\subsection{诊断}

侦听网按1000 m设计，偏保守：若真实接收半径接近1500 m，网仍然有效，只是偏密。9盘加连续残留的检测定理只在全向假设、$R_{\mathrm{eff}}\ge 1000$ m、$N\in[10,16]$且覆盖循环按前三类退出正常结束时保证听到。$A_{\max}=900$是算法自设闸，原题未规定；程序运行时间、测试窗口与技术异常同样是残差。覆盖路点静默扫描、射线归航与收尾追逐在这些情形返回时，尚未扫描的透镜可以仍与残留相交。同构几何核验动作次数至多213、演练132--160只说明这些样本没有碰到动作闸，不能把该闸从残差里删掉，也不证明覆盖循环必在有限步内走完。该定理不保证一次测向就能把$R_{\mathrm{sec}}$压到20 m，也不能搬到定向瓣上。有限角平分样本已空不是连续残留为空；同构几何核验与演练全部清除不是该定理。轴上弦长清除对一切$|\alpha|\le 1^\circ$保证命中当且仅当$r\le 20/(2\sin 0.5^\circ)$，覆盖对象是$\|SP\|=\|SG\|=r$时的弦圆；$r\tan 1^\circ$只是垂直截面半宽的保守充分条件。扩展搜索、横向有限振幅与沿示向度15 m不是这个圆。横向离散表含$\pm 34$ m，大于1500 m截面半宽26.18 m，因而26.18 m不是该表上界。射线归航后沿示向度再\texttt{/clear} 15 m发生在归航尚未清除且动作预算仍有剩余时，并不等到预算用尽，也不能回写成$R_{\mathrm{sec}}\le 20$ m的充分条件。问题二的交付点不是本问默认第二站。示向度误差在同一地点固定，重复在同一点检测不缩小楔形，必须换地点。覆盖目标走开放TSP，实际走法随源的分布而变，不能预先写成固定环序；共享D-最优站按式\eqref{eq:dfirst}在路程差不超过80 m时抢先作为下一站；该站落在覆盖路点120 m邻域内或覆盖剩余透镜、且路程差不超过220 m时同样抢先。覆盖未完时一站示向度并入当前路径的离原点条件是第一测向站到原点大于200 m。覆盖未完时两站服务点进入开放TSP不要求清除证书；选中后先圆心再测，仍开放则沿测得示向度归航。覆盖未完时失败回退仍是与残留相交的9网点，覆盖义务不改写成命中即射线的星形。现实时间用\texttt{/enter}返回的剩余秒数约束。虚拟时间可远长于现实时间。演练三次对应题目要的清除比例估计：每次以返回的总数为分母。演练结束时模拟器返回总数，因而比例可算；正式测试见表\ref{tab:official3}，只报告清除个数、平均定位清除时间与程序运行时间，不得把12、11、13写成干扰源总数。二者分表填写。原点单独覆盖失败是本策略存在的理由，而不是实现细节。

\section{问题四：混合源检测网与保扇形追逐}

\subsection{机理}

定向源在位置$G$、朝向$H$时，检测点$S$能收到信号当且仅当
\begin{equation}
\|S-G\|\le R_{\mathrm{eff}}\quad\text{且}\quad (S-G)\cdot u_H\ge 0,
\label{eq:dir}
\end{equation}
其中$u_H=(\cos H,\sin H)$。从$G$看，闭瓣是过$G$且法向为$u_H$的闭半平面。如图\ref{fig:sector}所示，有效覆盖是以$G$为顶点、朝向两侧各90度的闭扇形。最小接收半径只约束距离因子：落在背面的检测点即使$\|S-G\|\le 1000$ m也给出\texttt{no\_signal}。因此，单次无信号只能说明当前检测点未满足距离与前瓣的联合条件；问题三在正常结束时成立的全向检测结论适用于全向源。光学清除只依赖距离，问题一的圆心规则仍然充分。

本问只采用一套检测网。记$N_9$为问题三的原点加半径$1000\,\mathrm{m}$的8等分环，$R_{12}$为半径$2000\,\mathrm{m}$上的12等分点，检测网
\begin{equation}
\mathcal{N}=N_9\cup R_{12}.
\label{eq:q4net}
\end{equation}
全向源的距离可听性仍由$N_9$的闭式覆盖半径给出。定向源的可听性按式\eqref{eq:dir}逐点判定：当且仅当某个网点同时落在距离盘与前瓣内时，源才会被听到。该检测网对定向源的可听性取决于源的位置与朝向；当接收半径内的网点均位于开背瓣时，相应的位置--朝向组合构成残留集
\begin{equation}
\mathrm{Res}=\bigl\{(g,u):\mathcal{N}\cap\overline{\mathbb{D}}(g,1000)\ \text{全落在开背瓣}\bigr\}.
\label{eq:q4res}
\end{equation}
如图\ref{fig:dirlisten}所示，源在$G=(1200,0)$、朝向正东时，外环最近点$S^\star=(2000,0)$落入前瓣与$1000\,\mathrm{m}$盘的交，因而被听到；落在背瓣内的网点即使距离足够也听不到。如图\ref{fig:spoke}所示，场地边界中缝$G=1800\angle 15^\circ$到最近外环点$S^\star=2000\angle 0^\circ$的弦长约$534.17\,\mathrm{m}$，小于$1000\,\mathrm{m}$，故每个边界位置都有外环点落在距离盘内。

\begin{figure}[htbp]
\centering
\includegraphics[width=0.98\textwidth]{figures/directional_listen.pdf}
\caption{定向源的距离与朝向联合判定（$R_{\mathrm{eff}}=1000\,\mathrm{m}$）。左：完整检测网，圆点、方点分别为$N_9$、$R_{12}$，黑实线为场地边界。右：$G=(1200,0)\,\mathrm{m}$，朝向正东；红点$S^\star=(2000,0)\,\mathrm{m}$距源800 m且在前瓣内，可听。蓝点处于距离盘内的背瓣，其中$S_b$距源仅200 m，仍不可听；灰色空心点距离超限。}
\label{fig:dirlisten}
\end{figure}

\begin{figure}[htbp]
\centering
\includegraphics[width=0.55\textwidth]{figures/q4_directional_sector.png}
\caption{定向源的有效覆盖为朝向两侧各90度的闭扇形。$G$在圆心，$H$指向前向瓣。}
\label{fig:sector}
\end{figure}

\begin{figure}[htbp]
\centering
\includegraphics[width=0.82\textwidth]{figures/q4_silent_spoke.png}
\caption{外环距离引理的中缝剖面：$G$在场地边界$1800\angle 15^\circ$，$S^\star$在$2000\angle 0^\circ$，弦长$d\approx 534.17$ m，虚线圆是以$G$为心的$1000\,\mathrm{m}$盘。}
\label{fig:spoke}
\end{figure}

\subsection{模型}

侦听网$\mathcal{N}$由内圈9盘与外环12点并成，见表\ref{tab:q4pts}。其基本形式为
\begin{equation}
N_9=\bigl\{(0,0)\bigr\}\cup\bigl\{1000(\cos\tfrac{k\pi}{4},\sin\tfrac{k\pi}{4}):k=0,\ldots,7\bigr\},
\label{eq:q4n9}
\end{equation}
\begin{equation}
R_{12}=\bigl\{2000(\cos\tfrac{k\pi}{6},\sin\tfrac{k\pi}{6}):k=0,\ldots,11\bigr\}.
\label{eq:q4r12}
\end{equation}
如图\ref{fig:scouts}右幅所示，内圈点落在$1000\,\mathrm{m}$圆上，外环方点落在$2000\,\mathrm{m}$圆上。内圈对静默频道做连续残留扫描；外环只在占用尚未达到16且静默频道仍在时开圆访问。占用达到16后剩余静默频道为空，外环扫描即可停止。附录表\ref{tab:q4visit}列出全部候选点位，实际访问顺序由路径规划确定。混编覆盖循环设置动作上限$A_{\max}=900$。当预算耗尽、程序超时或发生技术异常时，尚未访问的网点仍可能对部分定向源具有可听性，这类情形与问题三检测定理的非正常结束情形相对应。

\textbf{引理（全向，9盘）。}假设源为全向，$G\in\overline{\mathbb{D}}(0,1800)$，$R_{\mathrm{eff}}\ge 1000\,\mathrm{m}$。完成$N_9$后必能听到该源，因为覆盖半径$\rho(N_9)\approx 956.05\,\mathrm{m}<1000\,\mathrm{m}$。证明见问题三。本引理针对全向源，定向源的可听性还需结合前瓣条件判断。

\textbf{引理（外环距离）。}假设$g\in\overline{\mathbb{D}}(0,1800)$，$S^\star$为$R_{12}$中方位最近的点。则方位差$\Delta\theta\le 15^\circ$。对固定$\Delta\theta$，距离
\begin{equation}
d(r)=\sqrt{2000^2+r^2-2\cdot 2000\cdot r\cos\Delta\theta}
\label{eq:q4dmid}
\end{equation}
在$r\in[0,1800]$上于边界取最大。取$\Delta\theta=15^\circ$、$r=1800\,\mathrm{m}$，得
\begin{equation}
d_{\mathrm{mid}}=\sqrt{2000^2+1800^2-2\cdot 2000\cdot 1800\cos 15^\circ}\approx 534.17\,\mathrm{m}<1000\,\mathrm{m}.
\label{eq:q4mid}
\end{equation}
故每个边界位置都有外环点落在距离盘内。该引理刻画距离条件，前瓣条件由定向源检测引理补充。

\textbf{引理（定向源检测条件）。}假设源为定向，$R_{\mathrm{eff}}\ge 1000\,\mathrm{m}$。源被听到当且仅当存在$S\in\mathcal{N}$满足$\|S-g\|\le R_{\mathrm{eff}}$且$(S-g)\cdot u\ge 0$。若最近外环点满足$(S^\star-g)\cdot u\ge 0$（含径向外指），则该源被听到；若$\mathcal{N}\cap\overline{\mathbb{D}}(g,1000)$全落在开背瓣，则该位置--朝向组合属于残留集。该引理给出给定位置与朝向下的判定条件，检测保证的范围由源的朝向共同决定。

\textbf{引理（占用16）。}假设干扰源个数$N\le 16$且频道互不相同。若已经听到16个互异频道，则剩余静默频道为空，外环巡回的覆盖任务随之结束。

\textbf{引理（光学方格清除）。}边长25 m的闭正方形，格心到顶点距离$25/\sqrt{2}\approx 17.68<20$。故若$G$落在该格内，在格心\texttt{/clear}必中。当$\Omega\subseteq\overline{\mathbb{D}}(C,R_{\mathrm{sec}})$且$R_{\mathrm{sec}}\le 80$ m时，在该盘上铺设此方格即可形成清除证书。如图\ref{fig:optical}所示，虚线为清除半径20 m，实线为启用该证书的80 m盘。

\textbf{引理（精确示向度下沿测向线越过源）。}假设第一站$S$落在开瓣$\{X:(X-G)\cdot u_H>0\}$内，且第二站沿精确方位$u=(G-S)/\|G-S\|$行走。则当路程$t<\|G-S\|$时仍在开瓣内；当且仅当$t>\|G-S\|$时落入互补开半平面，该站为\texttt{no\_signal}。该结论依赖精确方位；示向度含$\pm 1^\circ$误差时，行走线可能偏离$G$。

\textbf{引理（侧向位移）。}记$v$为测得示向度的单位左转法向，$P(t)=S+tv$。$P(t)$对一切实数$t$都留在闭瓣内，当且仅当$v\perp H$。否则逃逸阈值为
\[
|t|=\frac{(S-G)\cdot u_H}{|v\cdot u_H|},
\]
当$|t|$超过该阈值时，侧向路径离开闭瓣。由于$t$与$H$的相对方位事先未知，算法在收到\texttt{no\_signal}后切换侧向。保扇形$(r,t)$列表用于降低搜索时间，其作用是生成满足扇形约束的候选测站。

时间层按三步执行。第一，内圈把剩余$N_9$按连续残留剪枝，原点优先，随后按当前位置的方位做逆时针蛇形；每位姿扫完静默频道，\texttt{near}就地清除，\texttt{direction}仅当估计点距当前不超过280 m时暂停扫描并进入追逐。第二，若占用仍低于16且静默频道仍在，从最近外环点起按较近的转向开圆访问$R_{12}$，已落在某发现站80 m内的外环点跳过。第三，已听但尚未清除的频道进入服务旅行商路径。该阶段承担干扰源发现与路径组织，清除操作在后续服务层完成。

追逐采用保扇形追逐（sector-preserving chase）。其基本形式为：在已测到示向度的站$S$上取
\[
S'=S+r\,u_{\theta}+t\,v_{\theta},
\]
其中$r$取短沿迹、$t$取侧向；若$S'$返回\texttt{no\_signal}，立即翻转$180^\circ$侧向假设。图\ref{fig:secchase}给出了两种情形：沿示向度越过$G$的第二站进入背面，而取$(r,t)=(60,500)$ m时第二站仍位于同一闭扇形内。两站示向度差小于$8^\circ$时视为共线，并改取纯侧向站拉开交角。清除半径20 m与朝向无关。问题二的解析交付点$(800,700)$ m属于另一组设计参数，本问定向追逐采用上述保扇形候选。

场地点的内点计划为
\[
(r,t)\in\bigl\{(150,\pm 400),\,(60,\pm 500),\,(250,\pm 200),\,(100,0)\bigr\}\,\mathrm{m};
\]
其中前两项$(150,\pm 400)$ m的模长均为
\begin{equation}
\sqrt{150^2+400^2}\approx 427.2\,\mathrm{m}.
\label{eq:q4step}
\end{equation}
外点第一站则取更贴边的
\[
(r,t)\in\bigl\{(80,\pm 450),\,(0,\pm 550),\,(160,\pm 300),\,(250,0)\bigr\}\,\mathrm{m}.
\]
共线侧向取$r=0$、$t\in\{\pm 500,\pm 350,\pm 700\}$ m。问题二交付点的站距为$\sqrt{800^2+700^2}\approx 1063.0$ m，与式\eqref{eq:q4step}对应不同的设计参数。

混编追逐写成有限次循环，上限取10次，分款见表\ref{tab:q4chase}。每一次循环按当前示向度条数分支。定位多边形的半径判定与问题三共用同一套规则，定向瓣只改变圆心再测的可听含义。

\begin{enumerate}
\item 仅一站示向度：按上列$(r,t)$表作保扇形第二站，\texttt{no\_signal}立即换侧。若该表推进了新的测向记录，进入下一次循环。若全部候选都未能推进，沿测得示向度按路程
\[
\{120,250,400,600,850,1100\}\,\mathrm{m}
\]
前进：返回\texttt{near}则清除；返回\texttt{direction}则更新最后命中点并继续前进；返回\texttt{no\_signal}则先对最后命中点与当前点的中点执行\texttt{/clear}，失败后再清最后命中点。当候选测站均未获得新示向度时，沿示向度继续搜索。
\item 两站示向度差小于$8^\circ$：改取纯侧向站。若侧向推进得到交角不小于$8^\circ$的新示向度，进入下一次循环；否则沿示向度继续搜索。
\item 两站非共线：先做多边形清除判定。证书半径仍是$R_{\mathrm{sec}}\le 20$ m或对径圆且$D/2\le 20$ m。若$20\,\mathrm{m}<R_{\mathrm{sec}}\le 1000$ m，仍在最小包围圆圆心测量：返回\texttt{near}则清除；返回\texttt{direction}则把新楔并入$\Omega$并进入下一次循环。圆心再测对定向瓣不再是“必听到”：源可能落在背面。若圆心返回\texttt{no\_signal}且$R_{\mathrm{sec}}\le 80$ m，先在圆心强制\texttt{/clear}，失败后再用18 m六邻点
\[
C+18\bigl(\cos\tfrac{k\pi}{3},\sin\tfrac{k\pi}{3}\bigr),\qquad k=0,\ldots,5
\]
补清除；$18+20=38$ m的邻点仅覆盖圆心附近区域，半径80 m的不确定盘仍需通过新的测向信息处理。六邻点之后若仍未清除，则将该定位区域记为开放状态。若$R_{\mathrm{sec}}>1000$ m，则跳过圆心再测并直接记为开放。对于开放定位区域，对最后一站再次使用保扇形$(r,t)$表并将新楔并入$\Omega$；当短沿迹仍未将$\Omega$压缩至证书半径时，该步骤承担主要的局部化作用。若保扇形测站未产生新的示向度，再沿示向度继续搜索。
\end{enumerate}

循环达到10次且频道仍未清除时，先在当前最小包围圆圆心执行一次\texttt{/clear}；若$R_{\mathrm{sec}}\le 80$ m仍未清除，则铺设边长25 m的光学方格；随后沿最后一站示向度按$\{120,\ldots,1100\}$ m搜索。$R_{\mathrm{sec}}>20$ m时，圆心强制清除属于尝试性操作，清除保证仍需由相应的几何证书给出。若静默频道尚存，则继续访问剩余网点；若$R_{\mathrm{sec}}>20$ m，则继续设置保扇形测站，直至满足证书条件或达到预算上限。

\begin{table}[htbp]
\centering
\caption{问题四混编追逐：分款与局部化循环}
\label{tab:q4chase}
\begin{tabular}{>{\raggedright\arraybackslash}p{0.46\textwidth}>{\raggedright\arraybackslash}p{0.46\textwidth}}
\toprule
条件 & 动作 \\
\midrule
一站示向度 & 保扇形短沿迹加侧向；内点首项$(150,\pm 400)$ m，模长$427.2$ m；\texttt{no\_signal}立即换侧 \\
一站保扇形已推进 & 进入下一次循环 \\
一站保扇形未推进 & 沿示向度继续搜索，路程为$\{120,\ldots,1100\}$ m \\
两站示向度差小于$8^\circ$ & 纯侧向站拉开交角 \\
共线侧向未推进 & 沿示向度路程 \\
两站非共线且$R_{\mathrm{sec}}\le 20$ m，或对径圆覆盖且$D/2\le 20$ m & 在该圆心\texttt{/clear}（清除证书） \\
$20\,\mathrm{m}<R_{\mathrm{sec}}\le 1000$ m & 仍在最小包围圆圆心测量 \\
圆心再测返回\texttt{near} & 在圆心清除 \\
圆心再测返回\texttt{direction} & 新楔并入$\Omega$，重入循环 \\
圆心再测返回\texttt{no\_signal}且$R_{\mathrm{sec}}\le 80$ m & 先圆心强制\texttt{/clear}，再18 m六邻点，然后才把判定标成开放 \\
多边形判定仍开放（含$R_{\mathrm{sec}}>1000$ m，或再测与六邻点仍未清除） & 对最后一站再次使用保扇形测站 \\
保扇形仍未推进 & 沿示向度路程 \\
循环10次后仍未清除 & 圆心执行\texttt{/clear}；$R_{\mathrm{sec}}\le 80$ m时铺光学方格；随后沿示向度搜索 \\
\bottomrule
\end{tabular}
\end{table}

\begin{table}[htbp]
\centering
\caption{问题四检测网$N_9\cup R_{12}$的构成}
\label{tab:q4pts}
\begin{tabular}{lccc}
\toprule
量 & 取值 & 个数 & 说明 \\
\midrule
内圈 $N_9$ & 原点加1000 m 8等分 & 9 & 全向闭式界$\rho\approx 956.05$ m \\
外环 $R_{12}$ & 半径2000 m、12等分 & 12 & 中缝$\Delta\theta\le 15^\circ$ \\
检测网合计 & --- & 21 & 定向检测受朝向限制 \\
光学方格边长 / m & 25 & --- & 格心外接半径$17.68<20$ \\
\bottomrule
\end{tabular}
\end{table}

\begin{figure}[htbp]
\centering
\includegraphics[width=0.98\textwidth]{figures/q4_sector_chase.png}
\caption{保扇形第二站。朝向$H$指向$S_1$所在半平面。左：沿示向度越过源，落入背面；右：$(r,t)=(60,500)$ m留在同一闭扇形内。}
\label{fig:secchase}
\end{figure}

\begin{figure}[htbp]
\centering
\includegraphics[width=0.55\textwidth]{figures/q4_optical_grid.png}
\caption{光学方格清除证书：边长25 m的闭正方形外接半径约17.68 m，小于清除半径20 m；实线盘半径80 m是启用该证书的限度。}
\label{fig:optical}
\end{figure}

\subsection{估计与检验}

检测引理基于式\eqref{eq:q4n9}--\eqref{eq:q4mid}的闭式界，抽样结果用于数值核对。全向完成$N_9$当且仅当$\rho\approx 956.05$ m$<1000$ m。外环中缝$d_{\mathrm{mid}}\approx 534.17$ m$<1000$ m，每个边界位置都有外环点落在距离盘内。定向源被听到当且仅当某网点同时满足距离与前瓣条件；残留集由接收半径内网点均位于开背瓣的位置--朝向组合构成。光学方格外接半径$17.68$ m$<20$ m，因此在格心执行清除可覆盖整个方格。800个位置--朝向样本中799个被听到，1个属于开背瓣残留；该抽样结果反映所测样本的可听性，检测性能仍由位置与朝向共同决定。

在官方模拟器上连续进行三个独立案例的演练，均采用上述$N_9\cup R_{12}$检测网与保扇形追逐策略，见表\ref{tab:q4policy}。三次案例编码为94W6-GKZD-JQQR-WHHE、U3UH-8T3Y-3TSU-Q2A5、U9B5-CN4M-ZBEN-Y3JH，定向源个数分别为3、8、3，三次均全部清除。上述结果验证了该策略在所测案例中的清除效果，一般情况下的未检测区域仍由式\eqref{eq:q4res}刻画。

\subsection{结果}

由表\ref{tab:q4cov}可知，检测网由9盘与12个外环点组成，合计21点；外环半径2000 m，最近方位差至多$15^\circ$，边界中缝距离约534.17 m，小于1000 m。由表\ref{tab:q4cert}可知，光学方格边长25 m，格心外接半径约17.68 m，小于清除半径20 m。如图\ref{fig:scouts}所示，右幅展示了内圈圆点与外环方点在场地圆盘中的分布。如图\ref{fig:covmap}所示，左幅4000个均匀圆盘样本在9点网上全部满足距离条件；右幅800个位置--朝向样本中799个被$N_9\cup R_{12}$听到，1个属于开背瓣残留。图\ref{fig:dirlisten}与图\ref{fig:spoke}分别展示了距离条件与前瓣条件的联合判定。

\begin{table}[htbp]
\centering
\caption{问题四检测网$N_9\cup R_{12}$}
\label{tab:q4cov}
\begin{tabular}{lr}
\toprule
量 & 取值 \\
\midrule
内圈点数 & 9 \\
外环点数 & 12 \\
检测网合计 & 21 \\
外环半径 / m & 2000 \\
最近方位差上界 / 度 & 15 \\
边界中缝距离 $d_{\mathrm{mid}}$ / m & 534.1666881402358 \\
中缝是否小于1000 m & 1 \\
9盘覆盖半径 $\rho$ / m & 956.0510880488383 \\
\bottomrule
\end{tabular}
\end{table}

\begin{table}[htbp]
\centering
\caption{问题四光学证书与定向抽样结果}
\label{tab:q4cert}
\begin{tabular}{lr}
\toprule
量 & 取值 \\
\midrule
光学方格边长 / m & 25 \\
光学格心外接半径 / m & 17.677669529663685 \\
光学方格覆盖条件 & 满足 \\
位置--朝向样本数 & 800 \\
样本可听个数 & 799 \\
\bottomrule
\end{tabular}
\end{table}

\begin{table}[htbp]
\centering
\caption{问题四演练的策略参数。三次均使用检测网$N_9\cup R_{12}$与保扇形短沿迹。}
\label{tab:q4policy}
\scriptsize
\begin{tabular}{lcccc}
\toprule
组别 & 侦听网 & 点数 & 第二站 & 站距 / m \\
\midrule
演练三次 & $N_9\cup R_{12}$ & 21 & 内点首项$(150,\pm 400)$ m & 427.2 \\
\bottomrule
\end{tabular}
\end{table}

由表\ref{tab:q4reh}可知，本文检测网加保扇形在三个独立案例上的干扰源总数为16、12、11，全向源分别占13、4、8，定向源分别占3、8、3，清除个数与总数相同，比例均为1。动作次数为147、332、294，平均定位清除时间为267.641 s、617.168 s、718.879 s。由表\ref{tab:q4reht}可知，虚拟时间分别为4282.259 s、7406.016 s、7907.673 s，路程分别为17.2 km、27.3 km、30.9 km。三次合计19595.949 s、39个源，合并不配对均值为502.460 s/源，平均路程25.1 km。如图\ref{fig:q4reh}所示，真实个数与清除个数三次重合。第一次占用达到16，外环巡回可以提前停止；后两次占用低于16，需要完成外环扫描，相应增加了移动路程与搜索时间。三次全部清除表明该策略在上述案例中有效，检测网对其他位置与朝向的覆盖仍受开背瓣残留限制。

\begin{table}[htbp]
\centering
\caption{问题四检测网$N_9\cup R_{12}$的三次演练}
\label{tab:q4reh}
\scriptsize
\begin{tabular}{llrrrrrr}
\toprule
演练 & 案例编码 & 总数 & 全向 & 定向 & 清除 & 比例 & 动作 \\
\midrule
1 & 94W6-GKZD-JQQR-WHHE & 16 & 13 & 3 & 16 & 1.000 & 147 \\
2 & U3UH-8T3Y-3TSU-Q2A5 & 12 & 4 & 8 & 12 & 1.000 & 332 \\
3 & U9B5-CN4M-ZBEN-Y3JH & 11 & 8 & 3 & 11 & 1.000 & 294 \\
\bottomrule
\end{tabular}
\end{table}

\begin{table}[htbp]
\centering
\caption{问题四三次演练的时间与路程}
\label{tab:q4reht}
\begin{tabular}{lrrrr}
\toprule
演练 & 虚拟时间 / s & 平均定位清除时间 / s & 路程 / km & 动作次数 \\
\midrule
1 & 4282.259 & 267.641 & 17.2 & 147 \\
2 & 7406.016 & 617.168 & 27.3 & 332 \\
3 & 7907.673 & 718.879 & 30.9 & 294 \\
合并不配对 & 19595.949 & 502.460 & 25.1 & --- \\
\bottomrule
\end{tabular}
\end{table}

\begin{figure}[htbp]
\centering
\includegraphics[width=0.98\textwidth]{figures/q4_rehearsal_bars.png}
\caption{问题四检测网$N_9\cup R_{12}$的三次演练：真实个数与清除个数重合；右幅为平均定位清除时间。}
\label{fig:q4reh}
\end{figure}

问题四的三次正式测试结果见表\ref{tab:official4}。定向源朝向由模拟器生成，模型不预设其分布。清除个数与平均定位清除时间由HTTP接口导出。正式测试未提供干扰源总数，因此仅报告清除个数，不计算清除比例。

\begin{table}[htbp]
\centering
\caption{问题四正式测试结果}
\label{tab:official4}
\begin{tabular}{lcccc}
\toprule
测试 & 案例编码 & 清除干扰源个数 & 平均定位清除时间 / s & 程序运行时间 / s \\
\midrule
测试1 & SZ4E-B3GU-6F8M-VX4P & 10 & 699.461 & 1.166 \\
测试2 & 6P77-F2QX-E6GE-FB7Y & 16 & 418.219 & 1.199 \\
测试3 & D7EY-MTV2-7Q3S-U7XT & 10 & 746.149 & 1.423 \\
\bottomrule
\end{tabular}
\end{table}

由表\ref{tab:official4}可知，三次正式测试的案例编码分别为SZ4E-B3GU-6F8M-VX4P、6P77-F2QX-E6GE-FB7Y、D7EY-MTV2-7Q3S-U7XT。HTTP成功\texttt{/clear}次数分别为10、16、10，平均定位清除时间分别为699.461 s、418.219 s、746.149 s，程序运行时间分别为1.166 s、1.199 s、1.423 s。第二场清除16个、均时418.219 s：占用达到16时剩余静默频道为空，外环$R_{12}$可以提前停止，与检测网时间层一致。三次合计HTTP清除36次、虚拟时间21147.610 s，合并不配对587.434 s/已清。正式测试采用与演练相同的混合源追逐策略：内圈按$N_9$的连续残留进行蛇形扫描，收到\texttt{direction}且估计点距当前不超过280 m时暂停扫描并进入追逐；占用未达16时继续完成$R_{12}$外环扫描，随后按服务旅行商路径处理已听到但尚未清除的频道。

\FloatBarrier
\subsection{诊断}

定向检测引理在$R_{\mathrm{eff}}\ge 1000$ m且瓣为闭$180^\circ$时给出混合源的检测条件，对定向源的检测性能由位置与朝向共同决定。全向源在完成$N_9$扫描后必被听到；定向源被听到当且仅当某网点同时位于距离盘与前瓣内，未检测区域由开背瓣残留刻画。三次演练及同构几何核验均属于有限案例验证，其中同构几何核验的8个场景全部清除、均时580.7 s。保扇形$(r,t)$列表用于改善搜索效率；证书动作包括在占用低于16时完成检测网剩余网点的访问，以及在$R_{\mathrm{sec}}\le 20$ m时于圆心执行清除。侧向$S+tv$仅在$v\perp H$时对所有$t$保持在闭瓣内；沿测得示向度越过源的半平面翻转需要精确方位条件。光学方格在$R_{\mathrm{sec}}\le 80$ m时提供清除证书；对于更大的定位区，圆心强制\texttt{/clear}属于尝试性操作。全向检测定理适用于问题三的全向假设、9盘、连续残留剪枝以及覆盖循环按前三类退出正常结束的情形；$A_{\max}=900$为算法设定的动作上限，程序运行时间、测试窗口与技术异常构成额外终止情形。混编覆盖循环在相同情形下结束。800点样本中799个可听，为检测条件提供了数值核对。第一次演练占用16，因而可以提前停止外环扫描；后两次占用均低于16，完成外环扫描增加了移动路程与搜索时间。频道扫描优先选择测向机当前所在频道，以减少换台耗时。沿示向度的长跨步可能使第二站进入$180^\circ$瓣背面，因此定向源的第二站采用保扇形策略。演练结束时模拟器返回总数，据此计算清除比例；正式测试的清除个数、平均定位清除时间与程序运行时间见表\ref{tab:official4}。

\section{灵敏度分析}

\subsection{横向偏移}

灵敏度分析围绕影响设计决策的参数展开。由表\ref{tab:q2sens}可知，固定轴向距离1100 m时，横向偏移从200 m增到800 m，中位直径由93.12 m降到58.03 m，下降35.09 m，因此选取较大的横向偏移。$t$增至1000 m时，1000 m可听比例降到约0.07，直径中位数回升到60.29 m，候选带的$|t|$上限取约845 m。如图\ref{fig:sens}所示，该剖面为$(r,t)=(1100,800)$ m，与交付点$(800,700)$ m对应不同坐标。交付点$|t|=700$ m由$r=800$ m处轴上$\psi\ge 45^\circ$覆盖条件给出，即$|t|\ge 1500-800$；网格分位点$690.9$ m仅用于分位数设计，不能直接替代该解析条件。网格分位点远轴$\psi\approx 44.63^\circ$，低于该引理的角度阈值。

\subsection{楔半角}

楔半角扫描取剖面$(r,t)=(1100,500)$ m。半角从$0.5^\circ$增至$1^\circ$时，一次清除由1变为0；增至$1.5^\circ$时，中位直径升至92.65 m。如图\ref{fig:errsweep}所示，绿虚线在$1^\circ$处发生该变化，对应剖面为$(1100,500)$ m。网格分位点$(r,t)=(800,690.9)$ m的一次清除比例为$\eta_1=0.21$，中位直径为48.27 m，与上述剖面对应不同；交付点取$(800,700)$ m。题设$1^\circ$在$(1100,500)$ m处位于一次清除的临界位置。

\subsection{外环等分数}

问题四的外环距离条件随等分数$n$变化。半径固定为$2000\,\mathrm{m}$时，边界中缝距离为
\[
d(n)=\sqrt{2000^2+1800^2-2\cdot 2000\cdot 1800\cos(\pi/n)}.
\]
$d<1000$ m表示最近外环点落在距离盘内。如图\ref{fig:ablation}所示，$n=6$时$d\approx 1002.31$ m，超过1000 m；$n=12$时$d\approx 534.17$ m，落在盘内。由表\ref{tab:ablation}可知，$n=8$时$d\approx 766.86$ m已满足距离条件，增至16点后中缝降到约422.31 m，同时外环路程增加。综合距离条件与移动代价，交付方案取12点。该扫描保持检测模型不变，仅考察等分数对边界中缝距离的影响。定向听力仍受开背瓣限制，外环加密不会消除朝向因素。

\begin{figure}[htbp]
\centering
\includegraphics[width=0.82\textwidth]{figures/q4_ring_ablation.png}
\caption{外环边界中缝距离随等分数的变化。水平虚线为$1000$ m；星号为交付$n=12$，叉号为不合格$n=6$。}
\label{fig:ablation}
\end{figure}

\begin{table}[htbp]
\centering
\caption{外环等分数扫描：中缝距离是否小于1000 m}
\label{tab:ablation}
\begin{tabular}{rrrrr}
\toprule
$n$ & $\Delta\theta$ / 度 & $d_{\mathrm{mid}}$ / m & $d<1000$ & 外环半径 / m \\
\midrule
6 & 30 & 1002.31 & 0 & 2000 \\
8 & 22.5 & 766.86 & 1 & 2000 \\
10 & 18 & 626.41 & 1 & 2000 \\
12 & 15 & 534.17 & 1 & 2000 \\
16 & 11.25 & 422.31 & 1 & 2000 \\
\bottomrule
\end{tabular}
\end{table}

\subsection{场地裁剪}

问题一的场地裁剪把对径圆覆盖率从1降到0.955，表明直径圆的覆盖结论需要结合场地边界说明。问题三把侦听半径取1000 m，覆盖率为1，方案具有一定保守性并满足执行约束。

这些扫描均在同一模型框架下进行，楔形交、直径、侦听覆盖与追逐判定保持不变。侦听点加密会增加移动代价，正文报告满足覆盖约束的候选网络。

\section{模型评价}

\subsection{优点}

本文将示向度误差表示为凸楔形，将交会区域表示为半平面裁剪结果，并以最小包围圆与清除半径的关系判定清除条件。问题一给出对径圆覆盖条件，并以等边三角形与三站六边形说明场地构型的影响。问题二利用轴上$\psi\ge 45^\circ$覆盖条件确定第二站，交角中位数接近直角。问题三以9点侦听网的闭式覆盖半径和连续残留剪枝为基础，在覆盖循环正常结束时给出全向源的检测正确性；时间层先访问覆盖骨架，两站服务点在覆盖未完时进入开放TSP并在选中后进行圆心再测与归航，一站示向度通常在覆盖完成后沿测得方向搜索。算法动作上限取$A_{\max}=900$。问题四在9盘之外增设半径2000 m的12等分点，将定向源检测表示为距离条件与前瓣条件的联合判定，并使用保扇形第二站与光学方格完成局部化。主要定量结果包括：正交例直径34.92 m；网格分位点直径中位数48.27 m、一次清除比例$15/70$；交付点横向700 m满足轴上$\psi\ge 45^\circ$；9点覆盖半径956.05 m；光学方格外接半径17.68 m。论文体例遵循竞赛格式规范\cite{format}。

\subsection{不足}

第一，有效接收半径未知，侦听网按下限设计；当真实半径更大时，网络密度和移动路程相应增加。第二，题设$1^\circ$下两站一次清除比例为$\eta_1=15/70=0.21$，全向局部化中的覆盖骨架、D-最优插入与横向有限振幅属于时间启发式，虚拟时间对源的个数较为敏感。第三，定向检测的可听性受开背瓣影响，一次测向未必能将$R_{\mathrm{sec}}$压缩至20 m；光学方格在$R_{\mathrm{sec}}\le 80$ m时提供清除证书。占用数低于16时需要完成外环扫描，从而增加距离引理对应的移动代价。三次演练全部清除仅反映所测案例的结果，定向检测仍由开背瓣残留集描述。第四，侧向路径保持在闭瓣内的充要条件为$v\perp H$，保扇形$(r,t)$表用于改善搜索效率。第五，演练比例以模拟器返回的总数为分母，正式测试见表\ref{tab:official3}与表\ref{tab:official4}，分别报告清除个数、平均定位清除时间与程序运行时间。问题四检测网为$N_9\cup R_{12}$，定向源残留由开背瓣确定。连续残留检测定理适用于全向源、$R_{\mathrm{eff}}\ge 1000$ m、$N\in[10,16]$以及覆盖循环按前三类退出正常结束的情形；$A_{\max}=900$为算法动作上限，程序运行时间、测试窗口与技术异常构成额外终止情形。同构几何核验与演练结果用于经验验证。覆盖循环的有限步收敛性以及加密测站在有限步内达到$R_{\mathrm{sec}}\le 20$ m的性质，尚未在本文中证明。

\subsection{推广}

将楔形半角由仪器标定参数给出后，同一套裁剪与直径算法可用于其他测向平台。第二站的轴上覆盖点可根据清除半径重新计算；全向源采用覆盖骨架与后续示向度搜索，定向源采用短沿迹加侧向策略。外环半径与等分数可按实际$R_{\mathrm{eff}}$重新设计，使边界中缝保持在距离盘内。若允许配置第二只机器狗，正交基线可由两台设备同时占用，从而降低虚拟时间中的移动代价。推广时仍应保留示向度误差模型，并将正式测试与演练数据分别统计。

AI工具使用情况见同目录文件《AI工具使用详情》。

\begin{thebibliography}{99}
\bibitem{stansfield1947} Stansfield R G. Statistical theory of d.f. fixing[J]. Journal of the Institution of Electrical Engineers, Part IIIA: Radiocommunication, 1947, 94(15): 762-770.
\bibitem{torrieri1984} Torrieri D J. Statistical theory of passive location systems[J]. IEEE Transactions on Aerospace and Electronic Systems, 1984, 20(2): 183-198.
\bibitem{nardone1984} Nardone S C, Lindgren A G, Gong K F. Fundamental properties and performance of conventional bearings-only target motion analysis[J]. IEEE Transactions on Automatic Control, 1984, 29(9): 775-787.
\bibitem{foy1976} Foy W H. Position-location solutions by Taylor-series estimation[J]. IEEE Transactions on Aerospace and Electronic Systems, 1976, 12(2): 187-194.
\bibitem{jung1901} Jung H. \"{U}ber die kleinste Kugel, die eine r\"{a}umliche Figur einschlie{\ss}t[J]. Journal f\"{u}r die reine und angewandte Mathematik, 1901, 123: 241-257.
\bibitem{welzl1991} Welzl E. Smallest enclosing disks (balls and ellipsoids)[C]//Maurer H. New Results and New Trends in Computer Science. Lecture Notes in Computer Science, vol. 555. Berlin: Springer, 1991: 359-370.
\bibitem{sutherland1974} Sutherland I E, Hodgman G W. Reentrant polygon clipping[J]. Communications of the ACM, 1974, 17(1): 32-42.
\bibitem{toussaint1983} Toussaint G T. Solving geometric problems with the rotating calipers[C]//Proceedings of IEEE MELECON. 1983.
\bibitem{dogancay2005} Dogancay K. Bearings-only target localization using total least squares[J]. Signal Processing, 2005, 85(9): 1695-1710.
\bibitem{bishop2010} Bishop A N, Fidan B, Anderson B D O, Dogancay K, Pathirana P N. Optimality analysis of sensor-target geometries in passive localization[J]. Automatica, 2010, 46(3): 479-492.
\bibitem{clarke1964} Clarke G, Wright J W. Scheduling of vehicles from a central depot to a number of delivery points[J]. Operations Research, 1964, 12(4): 568-581.
\bibitem{format} 全国大学生数学建模竞赛组委会. 全国大学生数学建模竞赛论文格式规范[S].
\end{thebibliography}

\appendix
\section{支撑材料}

\subsection{问题一顶点与Monte Carlo样本}

表\ref{tab:q1vert}的四个顶点复述如下，供独立核验直径34.92 m。Monte Carlo请求400个样本，种子2026，全部有界，382个对径圆覆盖，覆盖率0.955；正文直方图与散点使用这400行，不是子样本。

\begin{table}[htbp]
\centering
\caption{正交交会例顶点（与表\ref{tab:q1vert}相同）}
\label{tab:q1vertapp}
\begin{tabular}{crr}
\toprule
编号 & $x$ / m & $y$ / m \\
\midrule
1 & 500.000 & 517.765 \\
2 & 482.550 & 499.695 \\
3 & 500.000 & 482.844 \\
4 & 517.450 & 499.695 \\
\bottomrule
\end{tabular}
\end{table}

\begin{table}[htbp]
\centering
\caption{问题一有界样本容量}
\label{tab:q1mcmeta}
\begin{tabular}{lr}
\toprule
量 & 取值 \\
\midrule
种子 & 2026 \\
请求样本数 & 400 \\
有界样本数 & 400 \\
对径圆覆盖个数 & 382 \\
覆盖率 & 0.955 \\
CSV行数 & 400 \\
\bottomrule
\end{tabular}
\end{table}


\subsection{侦听点坐标}

问题三点位见表\ref{tab:q3xy}。问题四检测网为$N_9\cup R_{12}$，由内圈9点和外环12点组成，见表\ref{tab:q4visit}。内圈按连续残留蛇形访问，外环采用开圆巡回，表中列出候选点位及其角色。

{\scriptsize
\begin{longtable}{crrl}
\caption{问题四检测网$N_9\cup R_{12}$点位清单（内圈蛇形、外环开圆）}
\label{tab:q4visit}\\
\toprule
编号 & $x$ / m & $y$ / m & 角色 \\
\midrule
\endfirsthead
\toprule
编号 & $x$ / m & $y$ / m & 角色 \\
\midrule
\endhead
1 & 0.0 & 0.0 & 原点 \\
2 & 1000.0 & 0.0 & 内圈9盘 \\
3 & 707.1 & 707.1 & 内圈9盘 \\
4 & 0.0 & 1000.0 & 内圈9盘 \\
5 & $-707.1$ & 707.1 & 内圈9盘 \\
6 & $-1000.0$ & 0.0 & 内圈9盘 \\
7 & $-707.1$ & $-707.1$ & 内圈9盘 \\
8 & 0.0 & $-1000.0$ & 内圈9盘 \\
9 & 707.1 & $-707.1$ & 内圈9盘 \\
10 & 2000.0 & 0.0 & 外环$R_{12}$ \\
11 & 1732.1 & 1000.0 & 外环$R_{12}$ \\
12 & 1000.0 & 1732.1 & 外环$R_{12}$ \\
13 & 0.0 & 2000.0 & 外环$R_{12}$ \\
14 & $-1000.0$ & 1732.1 & 外环$R_{12}$ \\
15 & $-1732.1$ & 1000.0 & 外环$R_{12}$ \\
16 & $-2000.0$ & 0.0 & 外环$R_{12}$ \\
17 & $-1732.1$ & $-1000.0$ & 外环$R_{12}$ \\
18 & $-1000.0$ & $-1732.1$ & 外环$R_{12}$ \\
19 & 0.0 & $-2000.0$ & 外环$R_{12}$ \\
20 & 1000.0 & $-1732.1$ & 外环$R_{12}$ \\
21 & 1732.1 & $-1000.0$ & 外环$R_{12}$ \\
\bottomrule
\end{longtable}
}

\subsection{实现说明}

机器狗标识取队号，场地标识取默认值；接口请求按串行方式发送，等待上一次响应后再发送下一次请求。附录所列源文件实现楔形交、直径、侦听网与追逐算法。表\ref{tab:q3reh}记录9盘加连续残留策略的三次演练，表\ref{tab:q4reh}记录检测网$N_9\cup R_{12}$加保扇形策略的三次演练，正式测试结果另列于表\ref{tab:official3}和表\ref{tab:official4}。

{\small
\begin{tabular}{ll}
\toprule
路径 & 功能 \\
\midrule
\texttt{wedge\_quad/experiment.py} & 楔形交、直径与第二站比较 \\
\texttt{wedge\_quad/geom.py} & 半平面裁剪与最小包围圆 \\
\texttt{wedge\_quad/sensitivity.py} & 横向偏移与半角扫描 \\
\texttt{bearing\_chase/experiment.py} & 侦听覆盖 \\
\texttt{bearing\_chase/searcher.py} & 问题三全向源与问题四混合源的搜索实现 \\
\texttt{bearing\_chase/geom.py} & 追逐几何例程 \\
\texttt{paper/plot\_paper.py} & 展示用图与表 \\
\bottomrule
\end{tabular}
}

\end{document}
